\documentclass[11pt]{article}
\usepackage[margin=1.1in]{geometry}
\usepackage{amsmath,amssymb,booktabs}
\usepackage{algorithm}
\usepackage{algpseudocode}
\usepackage[hidelinks]{hyperref}
\usepackage{color}
\newcommand{\benedikt}[1]{{\color{red} Benedikt: #1 }}
\newcommand{\Ftwo}{\mathbb{F}_2}
\newcommand{\Feight}{\mathbb{F}_{2^8}}
\newcommand{\Fbig}{\mathbb{F}_{2^{128}}}
\newcommand{\eq}{\mathrm{eq}}
\newcommand{\code}[1]{\texttt{#1}}

\newenvironment{optsummary}{%
  \begin{quote}\small
  \begin{description}
  \setlength{\itemsep}{1pt}\setlength{\parsep}{0pt}\setlength{\topsep}{2pt}
}{%
  \end{description}
  \end{quote}
}

\title{Yukon Optimizations of Flock's Prover\\
{\large measurements on Apple M1 Max, single-threaded; zerocheck at $2^{18}$,
PCS open at $2^{16}$ BLAKE3 compressions --- and, from
Section~\ref{sec:union}, the union protocol at $m{=}32$ ($2^{18}$
compressions), multi- and single-threaded}}
\author{branch \code{snark-fast-improvements}}
\date{August--September 2026}

\begin{document}
\maketitle

\begin{abstract}
Seven changes to the RS (additive-NTT) zerocheck path, each kept only after a
paired, order-alternating A/B measurement with a decisive sign test. Together
they reduce zerocheck round~1 from $1477\,$ms to $1205\,$ms ($-18.4\%$),
round~2 from $433$ to $312\,$ms ($-28\%$), the rounds-3+ tail from $455$ to
$307\,$ms ($-33\%$), and whole-proof time by $\approx11\%$
single-threaded. Three are ideas ported
from an independently optimized fork of this prover (the ``challenge'' tree),
re-derived and re-implemented rather than copied; three originated here. The
recurring themes: reductions in fields of characteristic two are
$\Ftwo$-linear and can be deferred almost arbitrarily; work that a circuit
fixes structurally can be skipped with a one-word guard; and on Apple silicon
the boundary between the vector and general register files costs more than
most instruction counts. A later section extends the record to the PCS open,
whose dominant pass loses its per-slot field multiply to the same
$\Ftwo$-linearity argument (\S\ref{sec:composedfold}). The final part
(Section~\ref{sec:union}) records the campaign after the prover moved
onto the repository's union protocol: which of the earlier wins
survived the change of statement, which died with it, and the twelve
that were landed on the new protocol --- among them the statistics
ladder in the open, the any-size Merkle leaf batcher, and the
promotion of the AG zerocheck --- which together took the $m{=}32$
prove from $854$ to $438$\,ms multi-threaded and from $5.25$ to
$2.91$\,s single-threaded.
\end{abstract}

\section*{Preliminaries}

\emph{Shared definitions and the benchmark shape. The numbered sections
group the work by prover phase, in pipeline order; each \emph{subsection}
is one verified optimization, presented with its mechanism and its paired
measurement.}

\paragraph{Notation.}
A few recurring terms, used across multiple sections below:
\begin{itemize}
\setlength{\itemsep}{1pt}
\item \textbf{witness builder} (``the builder''): the code that serializes
  the R1CS witness into the packed bit-buffers $z$, $a$, $b$.
\item \textbf{drain} (the gather drain): the loop that walks the per-row
  lookup tables to accumulate round~1's output messages; detailed in
  \S\ref{sec:twolane}.
\item \textbf{lincheck stripe}: the repacking of the witness $z$, indexed
  by inner/outer position, that the lincheck phase consumes; defined
  fully in \S\ref{sec:stripe}.
\end{itemize}


Fields. $\Feight = \Ftwo[x]/(x^8+x^4+x^3+x+1)$ and the GHASH field
$\Fbig = \Ftwo[X]/(X^{128}+X^7+X^2+X+1)$, with the ring embedding
$\varphi:\Feight\hookrightarrow\Fbig$. Write $\alpha := \varphi(x)$ and
$\gamma := X$.

The witness is $z\in\Ftwo^{2^m}$; at the benchmark shape $m=32$
($2^{18}$ BLAKE3 compressions, $k_{\log}=14$ bits per block). After the
univariate skip, round~1 of the zerocheck views the $m$ index bits as
\[
[\;\underbrace{\lambda}_{6\ \text{(univariate)}}\;|\;
 \underbrace{s}_{3\ \text{(small)}}\;|\;
 \underbrace{b}_{4\ \text{(medium)}}\;|\;
 \underbrace{t}_{m-13\ \text{(rest)}}\;],
\]
and must emit, for each of the $64$ points $\lambda$ of the skip domain, the
quadratic message $P^{AB}(\lambda)$ and the linear message $P^{C}(\lambda)$,
plus a pair of \emph{banks} (the ring-switch helper $\hat s_{v,C}$) that keep
the small parity dimension unfolded: the first small coordinate $s_0$ is
\emph{not} summed out of the C-side accumulation --- round 1 emits one
accumulator per value of $s_0$ (two banks) instead of a single combined
value, because the c-claim's PCS opening needs the two halves separately
($\hat s_{v,C}$'s canonical form recombines them with weights $1$ and
$\alpha$; the wire message is their plain sum).

The pinned values below are protocol design, common to both trees --- not
one of this document's optimizations. The optimization content is
\emph{which direction} a kernel exploits them; the two directions are
stated in math at the end of this section.
The protocol pins the small and medium challenges to \emph{friendly} values:
$r_{6+i} = \varphi(v_i)$ with $v=(\mathtt{0xF7},\mathtt{0x53},\mathtt{0xB5})$,
and $\beta_i = \gamma^{2^{i-1}}/(1+\gamma^{2^{i-1}})$. These make the eq
weights geometric:
\begin{equation}
\prod_{i=0}^{2}\eq(r_{6+i},K_i) \;=\; C_s\,\alpha^{K},
\qquad
\prod_{i=1}^{4}\eq(\beta_i,b_i) \;=\; \gamma^{b}/D,
\label{eq:friendly}
\end{equation}
with $K=\sum K_i 2^i$, $C_s=\prod_i(1+r_{6+i})=\varphi(\mathtt{0x1C})$, and
$D=\prod_i(1+\gamma^{2^{i-1}})$. The baseline kernels exploit
\eqref{eq:friendly} through lookup tables: the geometric weights are
enumerated once into byte-indexed convert tables and applied by gathers.
Three optimizations below exploit it in the opposite direction, replacing
tables by folds at the actual challenge values: the C-side banks are
recomputed as a true multilinear fold at the pinned challenges, with the
over-folded eq factors undone by a closed-form constant that exists only
because of \eqref{eq:friendly} (\S\ref{sec:stripe}); the prep kernel's
weight $x^K$ is factored so one piece stays table-driven while the $x^2$
piece is applied as a shift-and-conditional-XOR directly on the operand
(\S\ref{sec:pmull}); and the PCS open builds its composed fold table by
walking the monomials $X^r e$ with 127 exact shift-and-folds instead of 128
general multiplies (\S\ref{sec:composedfold} --- resting on $\gamma = X$
being the generator rather than on \eqref{eq:friendly} itself).

The two directions, in math. Both evaluate contractions against the same
geometric weight vector: with $w_K = C\,\alpha^K$ for $K < n$ and
bit-inputs $u \in \Ftwo^n$, the quantity is
$\langle w, u\rangle = \sum_{K<n} u_K\, w_K$. The \emph{table} direction
precomputes the entire linear map once,
\[
T[u] \;=\; \sum_{K<n} u_K\, w_K
\quad\text{for all } u\in\Ftwo^{n}
\qquad(2^{n}\ \text{entries, built by subset-sum doubling}),
\]
after which each of $N$ streamed inputs costs a single gather,
$\langle w, u^{(j)}\rangle = T[u^{(j)}]$: setup $O(2^{n})$, then per-input
cost independent of $n$ and no multiplies at all. The \emph{fold}
direction never materializes $T$; it applies the weights level by level
to the resident data itself,
\[
u \;\leftarrow\; u_{\mathrm{even}} \;+\; \alpha^{2^{i}}\, u_{\mathrm{odd}},
\qquad i = 0,\dots,\log_2 n - 1,
\]
costing $n-1$ multiplies per resident vector and touching no table
memory. Hence the rule separating them: the table amortizes when one
fixed map meets a stream of many inputs ($N \gg 2^{n}/n$, the
load-port-bound drains); the fold wins when there is essentially one
resident object --- an accumulator, a bank, or the table $T$ itself while
it is being constructed (subset-sum doubling \emph{is} a fold;
\S\ref{sec:composedfold} is that observation at $n=128$). Folds at
freshly \emph{sampled} challenges (rounds two and onward) get no such
shortcut --- there a fold is an ordinary multiply.

\paragraph{Measurement protocol.}
All numbers are single-threaded (ST) at the shape above unless marked 8T
(eight threads). Each verdict is a paired A/B: both arms built once, one
discarded warm-up each, then eight pairs with the \emph{order alternating
within each pair} (a fixed order is biased by thermal drift by roughly the
size of the effects measured here). Phase times are the minimum over the
$\sim$5 zerocheck invocations within one run; we call each phase's measured
time slice its \emph{bucket}. A change is kept only on a decisive sign test
(a nonparametric test of whether the paired differences all share the same
sign); two of the six wins were $8/8$ with disjoint ranges. Measure only on
AC power with the battery above $\sim60\%$: low battery caps frequencies,
and fast-charging a nearly empty battery is worse.

\section{Witness generation}

Everything downstream consumes the bit-packed witness; its witness builder
(defined in Preliminaries) is the first timed phase.

\subsection{Streaming full-word publication}
\label{sec:streamwit}

\begin{optsummary}
\item[Baseline:] OR sub-word fields into pre-zeroed storage: a $\approx$130\,MB$\times$3 memset plus a read-modify-write on every store.
\item[Improvement:] A sequential writer that carries a partial word in a register publishes complete u64s in order and never reads the destination.
\item[Mechanism:] Eliminates the $\approx$130\,MB$\times$3 memset and the per-store read-modify-write entirely.
\item[Result:] ST $87$--$90\to64.8$--$68.9\,$ms ($-24\%$, $3/3$ disjoint); 8T $-20\%$.
\end{optsummary}

The BLAKE3 witness builder fills three bit-packed buffers (z, a, b; one bit
per R1CS wire) per compression block. The incumbent wrote them by OR-ing
sub-word fields into pre-zeroed storage: a $\approx$130\,MB\,$\times$3 memset,
then a read-modify-write on every store. But the circuit's rows are
\emph{contiguous} through the block's useful bits, so a sequential writer
that carries a partial word in a register can publish complete u64s in
order and never read the destination --- the memset and the RMW tax both
disappear. The one out-of-order region (the output chaining value, 256-bit
aligned) is reserved and overwritten once the final state is known. BLAKE3
compresses each block with 56 rounds of its inner G function (the
add-rotate-xor mixing step); this fixed 56-G schedule is fully unrolled
by hand, with each round's state/message array indices written as fixed
literals rather than computed at runtime, so the register allocator sees
every round's actual data dependencies directly. Verified bit-identical
to the OR-based builder, on real and padding slots, for both values of
the prefix carry bit.
Measured, paired: ST $87$--$90 \to 64.8$--$68.9\,$ms ($-24\%$, $3/3$
disjoint); 8T $-20\%$. Idea from the challenge tree; their further
$\approx$$3$--$7\,$ms SIMD-lockstep stage was priced against its
$\approx$400 lines and not taken --- the full network was later built
anyway and measured null in-prove (\S\ref{sec:zpool}).

\subsection{The vectorized-packing port that became an allocation fix}
\label{sec:zpool}

\begin{optsummary}
\item[Baseline:] The lincheck stripe buffer was a fresh \emph{zeroed} $128$--$512$\,MB allocation on every prove; its first-touch page faults dominated the witness bucket.
\item[Improvement:] Pool that one buffer (an untyped byte pool alongside the field-element pool) instead of reallocating and zeroing it each prove.
\item[Mechanism:] Not an operation-count change but an allocation fix: removes the per-prove first-touch page-fault cost of a fresh $128$--$512$\,MB zeroed allocation, at $\approx$40 lines (versus the ported vectorized-packing network's $\approx$200 lines, which measured null).
\item[Result:] In-prove witness bucket $18.5\to8.0\,$ms at $n=2^{16}$ and $74\to33\,$ms at $n=2^{18}$, both $2/2$ paired at 8 threads.
\end{optsummary}

The challenge tree's last witness edge was believed to be its lane-wise
vectorized packing network: each bit-packed output stream owns a
u32-granular writer whose pending word lives in a NEON lane, every wire
bit is pushed by a single \code{vsli} (shift-left-insert) with
compile-time shift constants, four compression blocks proceed in
lockstep, and finished words dump contiguously through a
\code{vld4} de-interleave. We reproduced the full design with a
\code{build.rs} generator ($\approx$200 committed lines emitting the
unrolled kernel, rather than porting their $2.6$k generated lines). It
matched the existing witness builder's output bit-for-bit on its first
full test run --- and measured a \emph{null} in-prove at both shapes.

The discrepancy exposed the real cost, which was never the witness builder: the
lincheck stripe buffer --- a repacking of $z$ consumed by the lincheck
--- was a fresh \emph{zeroed} $128$--$512$\,MB allocation on every
prove, and its first-touch page faults dominated the bucket. Pooling
that one buffer (an untyped byte pool alongside the field-element pool)
took the in-prove witness bucket from $18.5$ to $8.0$\,ms at $n=2^{16}$
and $74$ to $33$\,ms at $n=2^{18}$, both $2/2$ paired at 8 threads ---
more than the packing port's entire predicted value, at $\approx$40
lines. The mechanism is thread-count-agnostic (the buffer is faulted
once per prove regardless); the 8-thread figures are simply where the
campaign measured it. The quad kernel was reverted unmerged; its design
survives in the commit history. This is the third time the campaign's
largest available win sat in the allocation story rather than the
kernel (the constant-region elision of \S\ref{sec:elision}, the pooled
round-1 precompute, and this): audit where the buffers come from before
porting compute.

\section{Zerocheck round 1}

Round~1 is the zerocheck's largest phase; the four kept changes below took it
from $1477$ to $1205$\,ms ($-18.4\%$) at the benchmark shape.

\subsection{Unreduced accumulation with multiplicative weight splitting}
\label{sec:pmull}

\begin{optsummary}
\item[Baseline:] Each $y_K$ computed \emph{reduced}: two \code{PMULL}s for the raw product plus a reduction costing four more, then a widening shift.
\item[Improvement:] Split the weight $x^K$ multiplicatively so the reduction is deferred to the final fold; accumulate unreduced.
\item[Mechanism:] $6$ \code{PMULL} $\to$ $2$ \code{PMULL} per row-block (Algorithm~1).
\item[Result:] Round~1 $1401.4\to1273.4\,$ms ($-9.1\%$), $8/8$, every pair in $[-8.7\%,-10.1\%]$.
\end{optsummary}

The prep kernel --- round~1's per-row operand-weighting stage --- computes,
per $K$-row ($K\in\{0,\dots,7\}$) and 16-lane block,
$y_K = a_K \odot b_K \in \Feight$ and accumulates $\sum_K x^K y_K$ in 16-bit
lanes, reducing once at the end. The baseline computed each $y_K$
\emph{reduced}: two \code{PMULL}s for the raw product plus a reduction costing
four more, then a widening shift by $K$ --- six \code{PMULL}s per row-block for
a reduction that the final fold makes redundant.

Let $u_K$ be the \emph{unreduced} 15-bit product. Then
$\sum_K x^K u_K \equiv \sum_K x^K y_K \pmod{p_8}$, but $x^K u_K$ overflows 16
bits for $K\ge2$. Split the weight so every factor lands where it is free:
\[
x^K \;=\;
\underbrace{(x^4)^{[K\ge4]}}_{\text{choice of gather table}}\cdot
\underbrace{(x^2)^{[\lfloor K/2\rfloor\ \mathrm{odd}]}}_{\text{byte-wise
doubling of the reduced operand}}\cdot
\underbrace{x^{K\bmod 2}}_{\text{in-lane shift}}.
\]
The $x^4$ factor costs nothing per element: a second copy of the table with
every byte pre-multiplied by $x^4$ scales each gathered XOR-sum by $x^4$, by
linearity of $T$. The $x^2$ factor is a six-instruction shift-and-conditional-XOR
on the already-reduced 8-bit operand ($\deg$ still $\le7$). The residual shift
is one vector instruction and raises term degree to at most $15$.

\begin{algorithm}[h]
\caption{Row-block accumulate, before and after}
\begin{algorithmic}[1]
\State \textbf{before:}\ $\mathit{acc} \mathrel{\oplus}=
  \mathrm{shift}_K\!\big(\mathrm{reduce}(\mathrm{pmull}(da, db))\big)$
  \Comment{$6$ \code{PMULL}}
\State \textbf{after:}\ \ $da \gets [K\ge4]\ ?\ \text{gathered from }x^4\text{
  table} : da$;\quad $da \gets [(K/2)\bmod 2]\ ?\ x^2\!\cdot\!da : da$
\State \phantom{\textbf{after:}}\ $\mathit{acc} \mathrel{\oplus}=
  \mathrm{shift}_{K\bmod2}\!\big(\mathrm{pmull}(da, db)\big)$
  \Comment{$2$ \code{PMULL}}
\end{algorithmic}
\end{algorithm}

Correctness needs the final reducer to be exact to degree $15$, one beyond its
documented ``$\le14$''; both the scalar and vector reducers were verified
exhaustively over all $2^{16}$ inputs (tests now permanent). Gather traffic is
unchanged apart from the second $16\,$KiB copy of the table.

\emph{Measured: round~1 $1401.4\to1273.4\,$ms ($-9.1\%$), $8/8$, every pair in
$[-8.7\%,-10.1\%]$. The challenge tree's attribution predicted
$100$--$140\,$ms; measured $128$.}

\subsection{The C-side message as a fold of the lincheck stripe}
\label{sec:stripe}

\begin{optsummary}
\item[Baseline:] The C-side banks, though a multilinear evaluation, were computed with the quadratic side's machinery: a bit transpose of every $8192$-bit window plus $32$ of the drain's $48$ gathers per lane.
\item[Improvement:] Compute the banks as a true multilinear fold of the already-materialized lincheck stripe at the pinned challenges, undoing the over-folded eq factors with a closed-form constant.
\item[Mechanism:] Removes the per-window bit transpose entirely; the drain's gathers per (lane, $b$) drop from three to one.
\item[Result:] Round~1 $1277.5\to1204.7\,$ms ($-5.7\%$), $8/8$; added-fold cost $\approx180\,$ms nets against $\approx250\,$ms of removals.
\end{optsummary}

In the fast hash setups $C=I$, so $\hat c = \hat z$ and everything round~1
emits on the C side is \emph{linear} in the witness. Concretely the two banks
are, for parity $p\in\{0,1\}$ and lane $\lambda$,
\begin{equation}
\mathrm{bank}_p[\lambda] \;=\;
\sum_{K\equiv p\ (2)} \alpha^{K}
\sum_{b} \frac{\gamma^{b}}{D}
\sum_{t} \eq(r_{\ge13}, t)\; z[\lambda,K,b,t],
\label{eq:banks}
\end{equation}
a multilinear evaluation --- which the baseline nevertheless computed with the
quadratic side's machinery: a bit transpose of every $8192$-bit window plus
$32$ of the drain's $48$ gathers per lane.

By \eqref{eq:friendly}, a \emph{true} multilinear fold at the actual pinned
challenges reproduces the weights in \eqref{eq:banks} exactly. The prover
already materializes the \emph{lincheck stripe}, a repacking of $z$ indexed
$[\,i_{\mathrm{inner}}\,|\,i_{\mathrm{outer}}\,]$, for the later lincheck
phase. So:

\begin{algorithmic}[1]
\State $v \gets \mathrm{fold\_stripe\_outer}(z_{\mathrm{stripe}},\
  \eq(r_{k_{\log}..m}))$
  \Comment{the one $O(2^m)$ pass; same kernel lincheck uses}
\For{$d = k_{\log}-1$ \textbf{downto} $7$}
  \Comment{data halves each round: $2^{14}\to2^{7}$ elements}
  \State $v[i] \gets v[i] + r_d\,(v[i]+v[i+2^{d}])$
\EndFor
\State $\mathrm{bank}_p[\lambda] \gets \dfrac{\eq(r_6,p)}{C_s}\;
       v[\,p\cdot64+\lambda\,]$
  \Comment{$C_s=(1{+}r_6)(1{+}r_7)(1{+}r_8)$, from $r$ itself}
\end{algorithmic}

\noindent
Dimension 6 (the parity) and dimensions $0..5$ (the lanes) are kept unfolded;
the constant in step~3 undoes the two eq factors the partial fold applied,
derived from the identity $Y_p = (C_s/\eq(r_6,p))\,\mathrm{bank}_p$ obtained
by factoring \eqref{eq:friendly} over the folded dimensions. With the banks
gone from the drain, the workers run AB-only: the per-window bit transpose is
deleted and the drain issues one gather per (lane, $b$) instead of three.
Equivalence was proven at two levels before wiring: the new banks
bit-identical to the drain's old output, and all three entry outputs
(dense and padded) identical between the two implementations.

\emph{Measured: round~1 $1277.5\to1204.7\,$ms ($-5.7\%$), $8/8$. The cost of
the added fold ($\approx180\,$ms) is included; the removals
($\approx250\,$ms) net out.}

\subsection{Instruction-level parallelism in the gather drain}\label{sec:twolane}

\begin{optsummary}
\item[Baseline:] One lane per iteration exposes three serial dependency chains (one per bank) to the out-of-order engine.
\item[Improvement:] Process two lanes per iteration so twice as many independent chains are in flight, with no change in total work.
\item[Mechanism:] 3 dependency chains in flight $\to$ 6 dependency chains in flight per iteration (same gather count and table size).
\item[Result:] Round~1 $-63.1\,$ms ($-4.3\%$), $8/8$; combined with \S\ref{sec:zeroskip}: $1477.3\to1403.1\,$ms ($-5.0\%$).
\end{optsummary}

The drain accumulates, per output lane, three XOR chains of depth
$n_{\mathrm{med}}=16$ (one per bank), each link a table gather feeding an XOR.
The gathers are independent but the chains are serial, so one lane exposes
three dependency chains to the out-of-order engine. Processing two lanes per
iteration doubles that to six with no change in total work:

\begin{algorithmic}[1]
\For{$\ell = 0,2,4,\dots,62$}
  \State $u_0,u_1,u_2 \gets 0$;\quad $v_0,v_1,v_2 \gets 0$
    \Comment{banks for lanes $\ell$ and $\ell{+}1$}
  \For{$b = 0 \dots n_{\mathrm{med}}-1$}
    \State $u_j \mathrel{\oplus}= \mathrm{convert}[b][\cdot]$;\quad
           $v_j \mathrel{\oplus}= \mathrm{convert}[b][\cdot]$
           \Comment{6 independent chains in flight}
  \EndFor
\EndFor
\end{algorithmic}

This targets latency, not throughput: an attempt to shrink the gather
\emph{table} ($64\to8\,$KiB, doubling gather count) had cost $3.6\%$ --- the
table already fit M1's $128\,$KiB L1D, establishing that the drain is bound by
gather count and chain depth, not footprint.

\emph{Measured: round~1 $-63.1\,$ms ($-4.3\%$), $8/8$. Combined with
Section~\ref{sec:zeroskip}: $1477.3\to1403.1\,$ms ($-5.0\%$).}

\subsection{Structurally fixed rows of the $b$ operand}\label{sec:zeroskip}

\begin{optsummary}
\item[Baseline:] Every $K$-row of $b$ is pushed through the $\Ftwo$-linear table $T$, including the $15$ of $256$ rows the circuit pins to $0^8$ regardless of input.
\item[Improvement:] Guard each row with a single zero check; since $T(0)=0$, a zero $b$-row contributes nothing and its work can be skipped outright.
\item[Mechanism:] Skips $64$ table loads and $4$ $\Feight$ multiplies for each of the $15$ structurally-zero $K$-rows (out of $256$).
\item[Result:] Round~1 $-42.4\,$ms ($-2.9\%$), $8/8$.
\end{optsummary}

Round~1's dominant kernel transforms the operands $a,b$ through a table $T$
implementing the inverse-NTT composition; $T$ is $\Ftwo$-linear, so
$T(0)=0$. An exhaustive tally of the packed BLAKE3 witness (256 word positions
per block $\times$ 256 blocks $\times$ 3 independent witnesses) shows the circuit
pins $38$ of $256$ eight-byte $K$-rows of $b$ regardless of the inputs, taking
only three values:
\[
\underbrace{\mathtt{0xff}^8}_{22\ \text{positions (const-one wires)}},\qquad
\underbrace{0^8}_{15\ \text{positions (structural zeros)}},\qquad
\text{one mixed word.}
\]
For a zero row the entire row's work vanishes: $db=T(0)=0$, hence
$y = da \odot db = 0$ contributes nothing to any accumulator. The kernel
gains a single guard:

\begin{algorithmic}[1]
\State \textbf{if} $\mathrm{load}_{64}(b\_\mathit{row}) = 0$ \textbf{then
return} \Comment{skips all 64 table loads and 4 $\Feight$ multiplies}
\end{algorithmic}

\noindent
No tally data ships and no positions are tracked: the guard is exact for any
witness, and a disagreeing witness falls through to the generic path.
(The all-ones rows were also tried: constant-folding them saves only the $b$
half of a row's work, and halving a row's loads halves its memory-level
parallelism, returning $6\,$ms where $27$ were predicted. Reverted; whole-row
elimination is what pays.)

\emph{Measured: round~1 $-42.4\,$ms ($-2.9\%$), $8/8$.}

\section{Zerocheck round 2}

Two kept changes, sharing the deferred-reduction machinery, took round~2 from
$433$ to $312$\,ms ($-28\%$).

\subsection{Deferred reduction in vector registers}\label{sec:wideneon}

\begin{optsummary}
\item[Baseline:] The $256$-bit accumulator lived as a four-word struct in general registers: six lane extracts per product to leave the vector file, four scalar XORs per accumulate.
\item[Improvement:] Keep the accumulator resident as two $128$-bit vector registers and form the unreduced product with Karatsuba.
\item[Mechanism:] Karatsuba product: three \code{PMULL}s instead of the schoolbook four; accumulate: four scalar XORs $\to$ two vector XORs; the six lane extracts move from per-product to once per worker chunk.
\item[Result:] Round-2 phase $433.1\to375.2\,$ms ($-13.4\%$).
\end{optsummary}

For $u\in\Ftwo^{256}$ (an unreduced carry-less product) let
$R(u)\in\Fbig$ be reduction mod $X^{128}+X^7+X^2+X+1$. $R$ is
$\Ftwo$-linear, so for any products $u_i = a_i \odot b_i$ (unreduced),
\[
\textstyle\sum_i R(u_i) \;=\; R\!\left(\sum_i u_i\right),
\]
and a sum of $n$ field products needs one reduction, not $n$. The baseline
already used this identity, but held the $256$-bit accumulator as a four-word
struct in \emph{general} registers: each product paid six lane extracts to
leave the vector file, and each accumulate was four scalar XORs.

The fix is representational, not algorithmic: keep the accumulator as two
$128$-bit vector registers. The unreduced product uses Karatsuba,
\[
(a_l + a_h X^{64})(b_l + b_h X^{64}) : \quad
\mathit{ll}=a_l b_l,\ \ \mathit{hh}=a_h b_h,\ \
\mathit{cross}=(a_l{+}a_h)(b_l{+}b_h)+\mathit{ll}+\mathit{hh},
\]
three \code{PMULL}s instead of the schoolbook four; accumulation is two vector
XORs; the four extracts are paid once per worker chunk on the way out, and the
existing scalar reducer is reused unchanged.

\emph{Measured: round-2 phase $433.1\to375.2\,$ms ($-13.4\%$).}

\subsection{The register-resident message loop}\label{sec:qres}

\begin{optsummary}
\item[Baseline:] The message loop crossed the vector/general register-file boundary three times per output pair (fold-accumulator extract, scalar field multiply in/out, re-entry for accumulation), at six \code{PMULL}s per pair.
\item[Improvement:] Keep the whole chain --- fold, in-register Karatsuba multiply, and accumulate --- resident in vector registers.
\item[Mechanism:] $6$ \code{PMULL} $\to$ $5$ \code{PMULL} per pair (three for the product, two for the fold through $X^{128}\equiv X^7{+}X^2{+}X{+}1$); all three per-pair register-file crossings removed, leaving only four stores and one eq load.
\item[Result:] Round-2 phase $358.1\to311.6\,$ms ($-13.0\%$), $8/8$, every pair in $[-12.6\%,-13.3\%]$.
\end{optsummary}

Round~2's message loop crossed the vector/general register-file boundary three
times per output pair: the table-driven fold extracted its vector accumulator
into a two-word struct; the message products
$g_1 = a_1 b_1$, $g_\infty = (a_0{+}a_1)(b_0{+}b_1)$ went through the scalar
field multiply (struct in, struct out); and the accumulate of
Section~\ref{sec:wideneon} re-entered the vector file. The rewrite keeps the
whole chain in vector registers: the fold returns its register unextracted;
the products use an in-register Karatsuba multiply (five \code{PMULL}s ---
three for the product, two for the fold through $X^{128}\equiv X^7{+}X^2{+}X{+}1$
--- against the baseline's six); the $\eq\cdot g$ products feed the
Section~\ref{sec:wideneon} accumulator directly. Per pair, the only remaining
memory traffic is the four required stores of folded values and one eq load.

\emph{Measured: round-2 phase $358.1\to311.6\,$ms ($-13.0\%$), $8/8$, every
pair in $[-12.6\%,-13.3\%]$, taken on a machine with an active browser
session (the min-of-5 protocol absorbs interactive bursts).}

\section{Zerocheck rounds 3+}

One structural change took the tail from $455$ to $307$\,ms ($-33\%$).

\subsection{Fusing away a memory pass}\label{sec:tailfuse}

\begin{optsummary}
\item[Baseline:] Two passes per worker chunk: fold and store $a,b$, then re-read the just-written multi-megabyte chunk from L2/DRAM to build the message.
\item[Improvement:] Fold each output pair in-register, store once, and consume the pair for the message while still in registers --- no re-read.
\item[Mechanism:] Memory traversals per worker chunk: $2\to1$; \code{PMULL} count is unchanged from the two-pass form.
\item[Result:] Rounds-3+ phase $449.5\to306.6\,$ms ($-31.8\%$), $8/8$, every pair in $[-31.1\%,-32.6\%]$ --- the largest single win of the effort.
\end{optsummary}

The tail rounds apply the register-resident treatment of \S\ref{sec:qres}
plus one thing round~2 had no opportunity for. The incumbent made \emph{two} passes per worker chunk:
fold $a_{\mathrm{in}},b_{\mathrm{in}}$ into $a_{\mathrm{out}},b_{\mathrm{out}}$,
then re-read the just-written multi-megabyte chunk to build the message ---
an L2/DRAM read of data that had been in vector registers moments earlier.
The fused kernel folds each output pair in-register
($e + r\,(e+o)$ via \code{mul\_q}), issues the one required store, and
consumes the pair for the message while still in registers:

\begin{algorithmic}[1]
\For{$x = 0 \dots \mathit{lo}-1$}
  \State $a_0,a_1,b_0,b_1 \gets$ fold four input pairs at $r$
    \Comment{in 128-bit NEON Q-registers}
  \State store $a_0,a_1,b_0,b_1$ \Comment{the required write, once}
  \State $g_1 \gets a_1 b_1$;\quad
         $g_\infty \gets (a_0{+}a_1)(b_0{+}b_1)$
         \Comment{\code{mul\_q}, never leaving registers}
  \State $\mathit{acc}_1 \mathrel{\oplus}= \eq_x \odot g_1$;\quad
         $\mathit{acc}_\infty \mathrel{\oplus}= \eq_x \odot g_\infty$
         \Comment{unreduced, \S\ref{sec:wideneon}}
\EndFor
\end{algorithmic}

\noindent
The \code{PMULL} count is identical to the two-pass form; the entire gain is
one traversal instead of two plus the removed struct crossings. Notably, two
earlier attempts on this exact loop had failed (a memory-interfaced pair-multiply
kernel, $+10\%$; wide accumulators alone, $0\%$): they changed the arithmetic
encoding and the accumulator while leaving the pass structure intact, which is
where the cost actually lived.

\emph{Measured: rounds-3+ phase $449.5\to306.6\,$ms ($-31.8\%$), $8/8$,
every pair in $[-31.1\%,-32.6\%]$ --- the largest single win of the effort.}

\subsection{The lookahead cascade: two rounds per pass}
\label{sec:cascade}

\begin{optsummary}
\item[Baseline:] After \S\ref{sec:tailfuse}, each tail round is one fused pass (fold at $\rho$ + next message): $\approx2L$ read $+\,L$ write per round, halving $L$ each time.
\item[Improvement:] Emit eight \emph{lookahead} partial sums per pass, answer the next round's message without touching the arrays, and materialize two folds in one $4{\to}1$ sweep --- round~2 hands the first lookahead to the tail, the exit hands the last fold to the boundary.
\item[Mechanism:] The post-fold message is a quadratic in the not-yet-drawn challenge whose coefficients are bilinear in pre-fold group values: eight per-group sums determine it for \emph{any} $\rho$. The enabling arithmetic (from the challenge tree's kernel): all eight products of a group share its eq weight, so pre-scaling the four $a$-rows once makes each product a single unreduced widening multiply --- $52$ vs $72$ PMULL per group. Intermediate half-size arrays are never written ($\approx55\%$ of two sequential rounds' traffic).
\item[Result:] $m{=}32$ 8T: tail $33$--$37$ vs $53$--$66\,$ms ($4/4$ disjoint); round 2 pays $+12$ (its integrated kernel also emits the sums --- the PMULL floor); zerocheck net $-8$--$13\,$ms (sign $3/4$). $m{=}30$: $-1.0$, no regression ($2/2$). Default on; \code{FLOCK\_NO\_ZC\_LOOKAHEAD=1} reverts. Transcript byte-identical by test.
\end{optsummary}

This reverses an early-campaign verdict: the same double-round idea lost
$7$--$15\%$ when first tried, and the regression was adjudicated to the
technique. It belonged to the kernel --- the scalar lookahead
accumulator spilled what the NEON one-weight-per-group form keeps in
registers. The technique and its implementation quality had to be judged
separately, and only the cross-tree anatomy exchange (their kernel
reached the same sums with 20 fewer PMULLs per group) exposed which was
which.

\subsection{Compact variant K: rounds 2--5 in two passes}
\label{sec:compactk}

\begin{optsummary}
\item[Baseline:] After \S\ref{sec:cascade}, round 2 ran as two sweeps (fold + store the full pairs, then an L2 re-read for the eight lookahead products), and the tail bound the challenges from materialized $64$-byte-per-pair state.
\item[Improvement:] The challenge tree's compact architecture, ported whole. Round 2's single sweep stores $48$ bytes per pair --- the folded even-row \emph{anchor} plus the packed-byte XOR of the pair, still in the bit domain --- and emits round 2's message AND round 3's as six deferred quadratic coefficients: the products are parity-split over pair index ($\mathrm{eq}_2(2y{+}1) = r_1\,\mathrm{eq}_3(y)$), so the odd halves ARE the round-3 aggregates after one $r_1^{-1}$. The consumer then binds $\rho_1$ and $\rho_2$ in one pass through two $\lambda$-scaled byte tables ($\mathrm{bound} = \mathrm{anchor} + \mathrm{fold}_\lambda(\delta)$ by fold linearity over $\mathbb{F}_2$; $\lambda_1 = \rho_1(1{+}\rho_2)$, $\lambda_3 = \rho_1\rho_2$), emits round 4's message, defers round 5's the same way, and materializes the quarter-size level where the cascade resumes unchanged.
\item[Mechanism:] Three stacked effects: one sweep instead of two (the parity split needs the same eight accumulators the lookahead used, but no re-read pass); $-25\%$ intermediate traffic in both directions; and challenge binding through table lookups instead of per-output multiplies. The degenerate-$b$ fast path rides along: an all-ones packed $b$ row folds to a per-table constant, so such pairs skip their 16 lookups and their provably-zero products --- value-preserving, and common in this circuit's $b$ operand.
\item[Result:] Producer $54.3$--$55.9\to48.0$--$51.0\,$ms ($3/3$); tail $31.9\to8.5$ with the new $\approx24$ ms double fold between them; mechanism sum $85.9$--$87.9\to79.9$--$84.5$ ($3/3$ disjoint, avg $-4.9$) at $m{=}32$ 8T grind-free. (An earlier parity claim compared against warm-window sub-line medians and is withdrawn: on bucket basis the zerocheck retains $+14$--$19\,$ms at 8T after this change.) Transcript-identical by three tests including a targeted $b\equiv1$ case.
\end{optsummary}

\subsection{Structured-$b$ shortcuts and the hetero round-1 drain}
\label{sec:structuredb}

\begin{optsummary}
\item[Baseline:] The round-1 prep kernel treated every 8-row block uniformly (its 16-bit geometric-eq machinery --- shared with the challenge tree since the round-1 parity port --- already skips zero $b$ rows); the round-1 drain ran on the P pool alone.
\item[Improvement:] Two exact runtime shortcuts dispatched on two integer compares per block: an all-ones $b$ block (the inv-NTT extension of the constant-one row is constant one, so $y_K = \mathrm{ntt}_a$) skips every $b$ gather and every product multiply; a single-live-$K_0$ block (rows with $b = 0$ contribute nothing by $\mathbb{F}_2$-linearity) collapses to one dual transform and one lane multiply. Separately, the drain's fold/reduce pulls chunks from a shared counter drained by the rayon pool and two utility-QoS E-core threads.
\item[Mechanism:] Both shortcuts are value-preserving algebra, not approximation, and both patterns are common in this circuit's $b$ operand; the challenge tree's further constants-census paths (positions hardcoded from a scored-distribution census) were left unported --- the SHA-256 control had already adjudicated that family as noise. Random test data never exercises either path, so the oracle test crafts $b$ with all-ones and single-$K_0$ quarters.
\item[Result:] Commit join wall (where the hoisted prep lives at 8T) $257.5$--$271.4\to252.5$--$259.7\,$ms, $5/6$ pairs, avg $-5.7$; drain hetero $-0.53\,$ms $3/3$ disjoint on its own line. Campaign-record run during certification: $454.70\,$ms / $576{,}521$ comp/s.
\end{optsummary}

\section{Lincheck}

\subsection{One word load per stripe, two stripes per pass}
\label{sec:lcfold}

\begin{optsummary}
\item[Baseline:] 16 loads per stripe step: 8 table entries plus 8 separate scalar byte loads for the indices.
\item[Improvement:] Load the 8 adjacent index bytes as one u64 with GPR shift-extracts, and fold two stripes per iteration so the table-entry XOR fuses to a single \code{EOR3}.
\item[Mechanism:] 16 loads/stripe $\to$ $\approx$9 loads/stripe.
\item[Result:] \code{partial\_fold\_z} ST $40.7$--$40.8\to27.8$--$28.0\,$ms ($-32\%$, $3/3$ disjoint); 8T $-28\%$.
\end{optsummary}

The lincheck's dominant pass folds the 128\,MB packed witness through
per-stripe 256-entry byte tables into eight Q-register accumulators. The
loop is load-port bound (M1: 3 loads/cycle), and the incumbent spent 16
loads per stripe step: 8 table entries \emph{plus 8 scalar byte loads} for
the indices. The indices are 8 adjacent bytes --- one u64 load and GPR
shift-extracts replace all eight --- and folding two stripes per iteration
lets the two table entries enter the accumulator as an XOR pair, which LLVM
fuses to a single three-input XOR instruction, \code{EOR3}, available under
ARM's SHA3 feature: $\approx$9 loads per stripe, same XOR multiset, output
bit-identical to the baseline. Measured, paired:
\code{partial\_fold\_z} ST $40.7$--$40.8 \to 27.8$--$28.0\,$ms ($-32\%$,
$3/3$ disjoint) --- landing on the pass's computed $\approx$28\,ms load
floor; 8T $-28\%$. This supersedes the earlier ``no reachable headroom''
verdict, which had priced blocking and table-geometry changes but not the
index-load restructure (idea from the challenge tree's asm kernel,
re-expressed in intrinsics).

\section{PCS open}

\subsection{Composing the block scalar into the fold table}
\label{sec:composedfold}

\begin{optsummary}
\item[Baseline:] The combine pass spent one field multiply plus one 16-load table fold per slot per claim --- $58\%$ of open time.
\item[Improvement:] Fold the block-constant multiply into a freshly composed per-(claim, block) table, built once from 127 shift-and-fold doublings plus subset-sum expansion, so the sweep only indexes the composed table.
\item[Mechanism:] Removes $2L\approx16.8$M per-slot field multiplies per open, replaced by a one-time $\approx$4.3k-word-op build per block ($127$ shift-and-fold doublings plus $16\cdot255$ subset-sum additions).
\item[Result:] Combine $94.4\to78.5\,$ms, open total $\approx160\to\approx145\,$ms ($-9\%$); at 8 threads combine is $10.9\,$ms, a near-linear transfer.
\end{optsummary}

\emph{Shape for this section: the $2^{16}$-compression open ($m=30$,
$L=2^{23}$ packed slots, two batched claims).} The ring-switch reduction
leaves, for each claim $k$, an eq tensor in factored form
$\eq = \mathrm{eq}_{\mathrm{lo}} \otimes \mathrm{eq}_{\mathrm{hi}}$
(lengths $B$ and $L/B$) and an $\Ftwo$-linear map
$L_k:\Fbig\to\Fbig$ (the $\gamma_k$-weighted byte-table fold). The
\emph{combine} pass materializes the batched basis
\[
b[j] \;=\; \sum_k L_k\bigl(\mathrm{eq}_{\mathrm{lo}}[j \bmod B]\cdot
                            \mathrm{eq}_{\mathrm{hi}}[\lfloor j/B\rfloor]\bigr),
\qquad j < L,
\]
and was the largest bucket of the open: one field multiply plus one
16-load table fold per slot per claim, $58\%$ of open time.

The multiply is removable. Multiplication by the block constant
$e=\mathrm{eq}_{\mathrm{hi}}[\lfloor j/B\rfloor]$ is itself $\Ftwo$-linear, so
the per-slot map $G_{k,e}(w) = L_k(w\cdot e)$ is a composition of
$\Ftwo$-linear maps and is fully determined by its $128$ monomial images
$g_r = L_k(X^r e)$. Once per (claim, block), encode $G_{k,e}$ as a fresh
$16\times256$ byte table: advance $X^r e$ by $127$ exact shift-and-fold
doublings (multiplication by $X$ costs one shift and a conditional XOR of the
reduction word), evaluate $L_k$ on each via the existing table, and expand
every byte position from its $8$ generators by subset-sum doubling
($16\cdot255$ additions). About $4.3$k word operations per block; the sweep
then indexes only the composed table, and the per-slot multiply --- $2L
\approx 16.8$M multiplies per open --- is gone.

The catch is the block length. The balanced split used by the many-pass
folds ($B = 2^{\lceil n/2\rceil} = 2^{11}$ here) is too fine for the build to
amortize; the deferred claims carry their own coarse split
$B = \min(2^{15},\,2^{\,n-8})$, making $\mathrm{eq}_{\mathrm{lo}}$ a single
$512$\,KiB L2-resident stream shared by every block and the build
$\approx0.4$ ops/slot. Correctness is exact, not approximate: every composed
entry is an XOR combination of exact field products --- the same XOR multiset
as the slot-multiply path --- so the transcript is bit-identical to the
slot-multiply path's (unit equivalence test; verification of full proofs
unchanged).

Measured (ST, minimum over samples): combine $94.4\to78.5\,$ms, open total
$\approx160\to\approx145\,$ms ($-9\%$); at $8$ threads the combine is
$10.9\,$ms, a near-linear transfer. The idea was found dormant in the
challenge tree (present, never isolated or measured there) and re-derived.

\paragraph{Rejected alternatives for the open.}
(i) Three-way XOR fusion (\code{EOR3}) in the fold's reduction tree is flat:
LLVM already emits it under \code{target-cpu=native}, and the sweep is bound
by its $32$ loads per slot, not its XOR count. (ii) Fusing both claims into
one sweep --- two live composed tables, one store, no read-back of the
intermediate --- \emph{regresses} $15\%$: $128$\,KiB of gather targets thrash
L1, so the claim-sequential order stands. (iii) The round-1 lookahead
coefficients accumulated in the combine's tail are an accounting wash, not a
speedup: tail-plus-initial-sumcheck costs $60.5$ vs.\ $60.6\,$ms across the
two trees' opposite placements of the same work.

\subsection{The direct open: sufficient statistics instead of a basis}
\label{sec:directopen}

\begin{optsummary}
\item[Baseline:] Each claim is materialized as a length-$L$ basis vector; the combine builds $b=\sum_k \gamma_k B_k$ and the sumcheck folds it beside the witness --- the basis' whole life is $\approx4L$ element-ops plus the matching DRAM traffic, and it is what made the open scale linearly in $L$.
\item[Improvement:] Never materialize the basis. Each claim carries a $128\times2^c$ banked statistic, captured for free where the witness already streams; the first $c$ sumcheck rounds run on the banks, and only the $L/2^c$ residual basis is ever built.
\item[Mechanism:] $B_k$ is a rank-1 eq tensor pushed through an $\Ftwo$-linear map --- the sumcheck messages it feeds are computable from a constant-size contraction, so the $O(L)$ representation was pure convenience.
\item[Result:] At $m=32$: open $97.4\to64.1\,$ms 8T ($3/3$ disjoint) and $614\to210\,$ms ST; whole proof $527$--$543\to484.3$--$518\,$ms, best $484.25\,$ms ($541$k c/s). Proof bytes identical; \code{FLOCK\_NO\_OPEN\_DIRECT=1} restores the basis path.
\end{optsummary}

\emph{Shape for this subsection: the ranked $2^{18}$-compression open
($m=32$, $L=2^{25}$, two batched claims).} After ring-switching, claim $k$
is the inner product $\langle f, B_k\rangle = t_k$ against
$B_k[i] = \varphi_k\!\left(\gamma_k\prod_j \eq(u^{(k)}_j, i_j)\right)$ ---
a rank-1 eq tensor at the claim point pushed elementwise through the
$\Ftwo$-linear ring-switch map $\varphi_k(u) = \sum_b \mathrm{bit}_b(u)\,
E_k[b]$. The incumbent evaluates this rank-1 object by writing out all
$2^\ell$ of its entries and then folding them $c$ times.

Split the word index $i = (e, h)$ into its $c$ lowest and $\ell-c$ highest
coordinates, so the tensor factors as $\mathrm{lo}_k(e)\cdot
\mathrm{hi}_k(h)$, and define the \emph{banked statistic}
\[
M_e^{(k)}[\beta] \;=\; \sum_h \mathrm{bit}_\beta\!\bigl(f[(e,h)]\bigr)\,
\mathrm{hi}_k(h)
\qquad (2^c \text{ banks of } 128 \text{ slices}),
\]
which is the ordinary $\hat s_v$ contraction with the low $c$ coordinates
left unfolded --- and which the prover computes \emph{anyway}: the AB
claim's banks fall out of lincheck's pre-sumcheck $z$-fold, the C claim's
out of zerocheck round 1's stripe fold, in both cases by stopping an
existing small fold $c$ coordinates early. Two states per claim then carry
the open: the per-bank transpose $A[b,e]$ (bits extracted \emph{before}
any $\Fbig$ binding --- the step that makes folding sound) and
$W[b,e] = \varphi_k(\gamma_k\,\mathrm{lo}_k(e)\cdot x^b)$. Both fold
bank-wise at the sumcheck challenges, each round's message is
$\sum_{e}\langle A(x_0,e), W(x_0,e)\rangle$ over the $128$-slice axis
--- $O(2^{c-r})$ inner products, no $O(L)$ touch --- and after the $c$
lane folds $W$'s surviving bank \emph{is} the byte-map generator from
which the $L/2^c$ residual basis is materialized and the incumbent
recursion resumes. With $c$ equal to the L0 fold count, the banked region
spans exactly the rounds where the basis was large. Every quantity is the
same field element the incumbent computes, reassociated --- so the proof
is byte-identical (test-pinned), and the kill switch restores the basis
path exactly.

Measured at the ranked shape (paired, alternating, same binary): open
$94.3$--$100.2\to59.5$--$72.1\,$ms at 8 threads ($3/3$ disjoint) and
$614\to210\,$ms single-threaded; zerocheck unregressed (the banked
captures ride existing folds); whole-proof best $484.25\,$ms / $541$k
compressions/s --- the campaign's best figure at this shape. This is the
table-vs-fold rule of the Preliminaries taken to its endpoint: the claim
is one resident rank-1 object, so contract once and never materialize the
map. Idea from the challenge tree's ranked open (located from its honest
per-phase numbers after an instrumentation artifact had hidden it);
algebra re-derived, implemented as four independently tested units.

\subsection{Refinements to the direct open}
\label{sec:directopenrefine}

\begin{optsummary}
\item[Baseline:] The freshly ported direct open spent $8.7\,$ms building its per-claim states (a 128-bit conditional scan per $W$ entry, serial per-bank transposes) and $26.4\,$ms on the L0 lane folds (one halving pass of $f$ per round, $\approx2L$ traffic).
\item[Improvement:] Build $W$ through the standard 16-lookup fold byte table; parallelize the bank transposes; and fold $f$ ONCE --- all $k$ lane challenges composed into a single $2^k{\to}1$ pass.
\item[Mechanism:] The composed fold is an option only the direct open has: its round messages come from the banked states and never touch $f$, so the witness is needed only at the level-1 boundary --- $\mathrm{out}[j] = \sum_{e<2^k}\mathrm{lag}[e]\,f[2^k j{+}e]$ with $\mathrm{lag} = \mathrm{build\_eq}(\rho_0..\rho_{k-1})$, the same field elements as $k$ successive folds at $\approx1.03L$ traffic instead of $\approx2L$.
\item[Result:] Open $59.7$--$61.3\to45.5$--$50.4\,$ms at 8T ($-20\%$, $3/3$ disjoint); whole proof $516$--$538\to492.5$--$511\,$ms. State construction $8.7\to1.4\,$ms (the challenge tree's floor for the same stage is $0.5$). Proof bytes identical throughout.
\end{optsummary}

Three local costs of the port, removed without touching its algebra. The
$W$-state entries $\varphi(\gamma\,\mathrm{lo}(e)\,x^b)$ are exactly
\code{fold\_b128\_elems}' per-element map, so the existing byte-table
build applies verbatim; the per-bank $128{\times}128$ bit transposes are
independent and parallelize; and --- the structural one --- because the
direct messages are computed from the banked states alone, nothing needs
the partially folded witness until the residual basis is materialized, so
the per-round fold chain collapses into one composed Lagrange pass. The
dense path could never do this: its every round message reads the folded
pair. Certified by a two-binary paired A/B (the three fixes share one
keep decision), transcript identity re-checked at each step.

\subsection{Truncating the induce transpose at rate $1/2$}
\label{sec:inducetrunc}

\begin{optsummary}
\item[Baseline:] The sparse-prefix induce scatter-transposes the query indicator over the full $2^{n+1}$ codeword domain, then \code{truncate}s to the low $2^n$ coefficients --- the final transpose layers compute a half that is thrown away.
\item[Improvement:] Fuse the final three transpose layers into one strided kernel that computes only the retained low half (idea from the challenge tree; exact precisely when $\log(1/\text{rate}) = 1$).
\item[Mechanism:] The transpose runs forward layers in reverse order, so the last three pair at distances $n/8$, $n/4$, $n/2$. An in-place 8-gather kernel applies the three butterfly levels in registers and writes the four low outputs; the final layer's retained output is the plain sum $a+b$, so its multiplies vanish entirely. Three full sweeps ($\approx6n$ traffic) become one $n$-read/$\tfrac{n}{2}$-write pass.
\item[Result:] Induce $6.8$--$7.6\to6.0$--$7.2\,$ms at $m{=}32$ 8T (trace attribution); certified jointly with \S\ref{sec:lazyood} (open bucket $40.8$--$41.3\to36.8$--$37.8\,$ms 8T, $-9\%$, grind-free protocol). Retained-half byte equality with the dense transpose by test.
\end{optsummary}

\subsection{Lazy out-of-domain binding}
\label{sec:lazyood}

\begin{optsummary}
\item[Baseline:] Each OOD sample $z$ builds the dense $\mathrm{eq}(z,\cdot)$ table, runs the fused message+eval pass over $(f, \mathrm{eq}_z)$, and combines $\alpha\cdot\mathrm{eq}_z$ into the running basis --- $\approx7L$ traffic per sample, five samples at L0 alone.
\item[Improvement:] Never materialize $\mathrm{eq}_z$: factorize it as $\mathrm{eq}_{\mathrm{lo}}\otimes\mathrm{eq}_{\mathrm{hi}}$, read only $f$ for the message, and defer every sample's basis combine into the level's basis glue, drained fused in its single pass (idea from the challenge tree).
\item[Mechanism:] The eq table is LSB-first, so splitting $z$ splits the table as a tensor: $\mathrm{eq}_z[x] = \mathrm{eq}_{\mathrm{lo}}[x \bmod 2^k]\cdot\mathrm{eq}_{\mathrm{hi}}[\lfloor x/2^k\rfloor]$. The hi factor is constant per $2^k$-block, so block partials against the L1-resident lo table are scaled once per block --- the same multiply count per pair as the dense kernel with no table build or re-read. Glue stashes $(\alpha\cdot\mathrm{eq}_{\mathrm{hi}}, \mathrm{eq}_{\mathrm{lo}})$; the next basis glue adds all pending tensors alongside $\beta\cdot b_{\mathrm{new}}$ in its one read-modify-write (fold paths flush as insurance). Bit-identical by distributivity and order-free $\mathbb{F}_{2^{128}}$ sums.
\item[Result:] OOD stage $2.6$--$3.0\to0.4$--$1.0\,$ms at $m{=}32$ 8T; the draining basis glue pays $+0.1$--$0.6$. Certified jointly with \S\ref{sec:inducetrunc}.
\end{optsummary}

\subsection{Fusing the boundary pass}
\label{sec:boundaryjoin}

\begin{optsummary}
\item[Baseline:] At the direct open's level-1 boundary, the composed witness fold (one streaming pass over the full witness, DRAM-bandwidth-bound) and the residual-basis build (byte-table folds, 16 L1 gathers per element) ran back-to-back: $4.5 + 5.6\,$ms.
\item[Improvement:] Two stages. First, \code{rayon::join} — they are data-independent and stress different resources, so they share the pool ($10.1\to7.3$--$8.4$). Then full fusion (from the challenge tree's materialize pass): one sweep over the witness produces \emph{both} arrays — the gathers and the factorized eq products execute in the witness stream's stall slots.
\item[Mechanism:] Pool-level overlap recovers only part of the idle issue slots; instruction-level fusion recovers the rest. Per output: the $2^k{\to}1$ Lagrange fold accumulates unreduced (3 vs 6 PMULL per product, one reduction), and $b_1[j] = \sum_b \mathrm{fold}_b(\mathrm{eq}_{\mathrm{lo}}[j \bmod B]\cdot\mathrm{eq}_{\mathrm{hi}}[j/B])$ — the suffix eq tensor is never materialized. The standalone split-tensor fold was adjudicated \emph{slower} (its per-slot multiply outweighed the saved tensor build as its own pass); fused under the stream, that verdict inverts. Exact: deferred reduction is linear over XOR, field multiplies are exact, sums are order-free.
\item[Result:] Boundary pair $10.1\to5.5$--$5.7\,$ms; the open's initial-sumcheck stage $11.3\to5.9$--$6.0\,$ms at $m{=}32$ 8T (challenge tree: $5.19$). Join stage certified with \S\ref{sec:blockedtranspose} (open sign $3/3$); fusion stage phase-paired $3/3$ disjoint ($5.85/6.04/6.59$ vs $7.52/8.40/8.13$) — the bucket-level A/B could not resolve the $\approx2\,$ms effect in that window's noise, the per-phase pairing could.
\end{optsummary}

\subsection{Cache-blocking the induce transpose}
\label{sec:blockedtranspose}

\begin{optsummary}
\item[Baseline:] The sparse-prefix induce's dense remainder ran one full-array parallel sweep per transpose layer --- ten sweeps over the $2^{21}$-element domain at $m{=}32$ L0, $\approx640\,$MB of traffic.
\item[Improvement:] Run every layer whose blocks fit an L2-resident chunk in a single parallel pass over chunks, layers applied back-to-back in cache.
\item[Mechanism:] The transpose processes high layers (small blocks) first, and blocks nest: a chunk sized to the run's lowest layer contains complete blocks of every higher layer, so a $2\,$MB chunk hosts nine consecutive layers with one read and one write of DRAM traffic ($\approx130\,$MB total) --- the same sub-group decomposition the interleaved encode NTT already uses for its deep layers, applied to the transpose direction.
\item[Result:] Induce $6.6$--$8.5\to5.4$--$6.7\,$ms at $m{=}32$ 8T. Certified jointly with \S\ref{sec:boundaryjoin}: open bucket paired sign $3/3$ ($-2.8/-3.2/-8.0$, avg $-4.7\,$ms), grind-free protocol. Byte-equality tests cover the blocked run with and without the fused low-half tail.
\end{optsummary}

\subsection{Recursive commits from the message}
\label{sec:recursivefill}

\begin{optsummary}
\item[Baseline:] Each recursive Ligerito level's commit replicated its folded witness into the codeword buffer (a full write plus read-for-ownership), then transformed from the rate layers: L1 alone is a 32\,MB fill ahead of the NTT.
\item[Improvement:] The first fused top pass reads its rows straight from the compact folded witness (\code{from\_message}), so the standalone replicate pass and its RFO traffic disappear.
\item[Mechanism:] The same fused pass had been built for the \emph{main} commit and adjudicated null there (its join window is compute-limited, so deleted traffic bought nothing); the recursive commits run standalone and bandwidth-bound, where the challenge tree's equivalent measures fill $= 0.00$. The reverted patch was resurrected for this path only.
\item[Result:] L1 fill $0.5\to0$, all-in NTT $5.2$--$6.9\to5.0$--$5.2$; recursive commits $14.3$--$15.5\to12.9$--$13.3\,$ms at $m{=}32$ 8T. Open bucket paired sign $3/3$ ($-2.35/-1.60/-0.56$, avg $-1.5$), totals $3/3$.
\end{optsummary}

\section{Merkle hashing: BLAKE3}

\subsection{Two transposed four-wide states in flight}
\label{sec:blake3neon}

\begin{optsummary}
\item[Baseline:] The crate's NEON backend runs a single 4-wide transposed state (one \code{uint32x4} per state word); BLAKE3's serial G function leaves Apple's four NEON pipes half idle.
\item[Improvement:] Interleave a second independent transposed 4-wide state, G-for-G, to fill the idle pipe slots.
\item[Mechanism:] 1 transposed 4-wide state (4 messages in lockstep) $\to$ 2 interleaved 4-wide states (8 messages in flight, the full 32-register NEON file); no added work per message.
\item[Result:] ST tree build $1.63\to2.30\,$GB/s at 512\,B leaves ($+41\%$) and $1.71\to2.42\,$GB/s at 1\,KiB leaves ($+42\%$).
\end{optsummary}

The PCS Merkle tree can hash with SHA-256 (hardware silicon, $\approx$2.5\,GB/s
per core) or BLAKE3. The blake3 crate's NEON backend processes four messages
in lockstep --- one \code{uint32x4} per state word --- but a single 4-wide
state is \emph{latency}-bound: BLAKE3's G function is a serial
add--xor--rotate chain, and Apple's four NEON pipes sit half idle waiting on
it. Interleaving a \emph{second} independent transposed 4-wide state, G-for-G
(eight messages in flight, the full 32-register NEON file), fills the pipes
--- the same independent-chains pattern as \S\ref{sec:twolane}. Dispatch:
groups of eight leaves or parent pairs take the new kernel; tails fall to the
crate path; byte-identical to the crate's original kernel by an equivalence
test over all batched message sizes, counts, and flag combinations.

Measured (ST tree build): $1.63 \to 2.30$\,GB/s at 512\,B leaves ($+41\%$,
now $1.08\times$ \emph{faster} than SHA-256) and $1.71 \to 2.42$\,GB/s at the
1\,KiB production-geometry leaves ($+42\%$, parity with the SHA silicon). The
challenge tree reaches $\approx$2.58\,GB/s with a 2.6k-line generated
assembly kernel; this is $\approx$290 lines of intrinsics within $6\%$ of it
--- LLVM absorbed the full register pressure without measurable spill cost.
End-to-end the change erases the $-5\%$ penalty BLAKE3-Merkle previously
carried in this tree, making the Merkle hash choice free.

\section{Multithreading}
\label{sec:mt}

The campaign target mirrors the single-threaded closure: CPU-only
multithreaded throughput within $10\%$ of the challenge tree. Start:
$1.19\times$ apart at the $2^{16}$ shape. Standing after the two
subsections below: $1.146\times$ (149.3 vs 130.2\,ms, paired same-minute,
production configuration). All numbers here are 8-thread-class production
runs at $2^{16}$ compressions, on Apple silicon's split between
performance cores (P-cores) and efficiency cores (E-cores). Several
results below are a paired \emph{kill-switch A/B}: the optimization is
gated behind a runtime toggle, so disabling it reproduces the pre-change
baseline in the same binary.

\subsection{Round-1 AB prep under the commit, on all cores}
\label{sec:abhoist}

\begin{optsummary}
\item[Baseline:] Overlapping the witness-only AB transform with the commit under \code{rayon::join} measured a wash twice on the P-core-only pool; the transform's output buffer was a fresh zeroed allocation.
\item[Improvement:] Draw the output buffer from the scratch pool uninitialized, and run the join window on the all-core (P+E) pool instead of P-cores only.
\item[Mechanism:] Not an operation-count change but a memory-traffic and scheduling fix: removes the memset-plus-page-fault cost of a fresh zeroed allocation, and moves the overlap window onto the otherwise-idle E-core issue capacity left by the bandwidth-bound commit NTT.
\item[Result:] Kill-switch A/B $3/3$ at $2^{16}$ ($149.6$--$151.3$ vs $152.5$--$156.3\,$ms), $2/2$ at the ranked shape (clean pair $-57\,$ms).
\end{optsummary}

Round 1's dominant half --- the shift-reduce AB transform --- reads only the
witness, so it can run before the commitment root exists. Overlapping it
with the commit under \code{rayon::join} was measured a wash TWICE on the
8-P-core pool: both arms time-slice a saturated pool. Two changes flip it
decisive: (i) the transform's $2^{m-13}$\,KiB output buffer comes from the
scratch pool \emph{uninitialized} (a fresh zeroed allocation's memset +
page faults consumed the entire overlap gain; slots for padding-skipped
rows stay stale-but-unread, bounded by the same padding table every
consumer uses), and (ii) the join window runs on an all-core (P+E) pool
while the rest of the prove stays on P-cores. The efficiency cores are the
active ingredient: the commit's NTT is bandwidth-bound and leaves issue
capacity, and the prep is gather/PMULL compute the E-cluster can add
without competing for DRAM. (The converse placement --- a bandwidth-bound
task like the lincheck stripe transpose on the E-cores during the commit
--- measured \emph{negative}.) Paired kill-switch A/B: $3/3$ at $2^{16}$
($149.6$--$151.3$ vs $152.5$--$156.3$\,ms), $2/2$ at the ranked shape
(clean pair $-57$\,ms). Bit-identical: the precomputed and inline
transforms are equality-tested on all outputs.

\subsection{The streaming commit: pipelined leaves and the fill's overlap budget}
\label{sec:streamcommit}

\begin{optsummary}
\item[Baseline:] The commit ran as stages --- replicate-fill, NTT, Merkle --- joined with the round-1 AB prep on the all-core pool. The window was compute-additive: the prep overlapped only the short fill stage, then $\approx180\,$ms of transform and hashing ran with the prep already finished ($\approx273\,$ms wall).
\item[Improvement:] The challenge tree's architecture, ported whole: the deep NTT pass publishes each finished sub-group as $\approx$1-MiB leaf jobs into a small bounded queue drained by utility-QoS helper threads (which the scheduler steers to the efficiency cores); each job hashes its leaves \emph{and} its aligned local parent subtree while the lines are hot; a full queue hashes inline on the publisher; the tail drains on the main pool; only the tiny shared top levels remain after ($0.03\,$ms). Crucially, the fill is fused into the first top pass (from-message), so the prep overlaps the \emph{entire} encode+leaves window.
\item[Mechanism:] The fill fusion is the load-bearing piece, not the pipeline: without it the prep burns its overlap budget beside a standalone bandwidth stage that needs none of it (the v1 port measured \emph{worse} than staged). The staged-shape adjudication of fill fusion --- ``deletes a read, not compute'', null --- does not transfer: in the pipeline shape the deleted stage is what buys the prep its hiding window. Null verdicts are architecture-scoped.
\item[Result:] Join wall (= commit + prep, one window) $260.5$--$262.6\to252.5$--$254.7\,$ms ($3/3$ disjoint); whole proof $479.8$--$502.8\to456.9$--$482.5\,$ms ($3/3$ disjoint, avg $-23$). Campaign record $456.95\,$ms / $573{,}681$ comp/s at $m{=}32$, grind-free. Tree and root bit-identical by test.
\end{optsummary}

\subsection{The ranked radix-8 top: dual from-message pass, staged stores, hetero tiles}
\label{sec:rankedtop}

\begin{optsummary}
\item[Baseline:] After \S\ref{sec:streamcommit}, the commit's top layers ran as a from-message fused-2 pass plus three fused-2 sweeps, all on the P pool (the E-core helpers idle until the deep pass publishes leaf jobs).
\item[Improvement:] Layers 1--9 as \emph{three} radix-8 passes: layer 1 is a dual-destination from-message kernel (one witness read produces both replica blocks; block 0's spine twiddles are zero, so its chain is XOR-only), layers 4 and 7 are in-place sweeps with q-register-resident butterflies; every pass distributes 128-row tiles across the rayon pool AND two utility-QoS E-core threads via one shared work counter.
\item[Mechanism:] Radix-8 was not the earlier fused-3 failure --- the kernel shape was. Fresh-destination scatter stores (sixteen 16-B streams, tens of MB apart) defeat the streaming-store detector and pay a read-for-ownership on the whole 1-GB destination; staging each row group in an 8-KB L1 tile and emitting sequential non-temporal bursts restores fill-like store behavior. In-place passes need no staging (their lines are already owned). One fewer full sweep of traffic, one witness read instead of two, and the E-cores now assist the top before switching to leaf hashing.
\item[Result:] Commit join wall $273.1$--$275.5\to256.6$--$266.7\,$ms ($3/3$ disjoint, avg $-13.1$; commit bucket $3/3$). Bit-identical by the from-message and commit equality tests.
\end{optsummary}

\subsection{Witness constant-region elision via scratch provenance}
\label{sec:elision}
\benedikt{what is this title? Use plain english please}
\begin{optsummary}
\item[Baseline:] The full-write witness builder rewrites every byte of z/a/b each prove, including regions --- $b$'s all-ones prefix, reserved-slot words, all three streams' zero tails --- identical across every prove of a layout.
\item[Improvement:] Tag a pooled buffer with its provenance --- which layout's output it currently holds, derived independently at give and take sites --- and, on a tagged hit, skip the constant-region writes.
\item[Mechanism:] No explicit byte or operation count is given in this section; this is a work-removal via provenance tracking, not a stated op-count reduction.
\item[Result:] Paired kill-switch A/B: $4/5$ ($147.0$--$149.0$ vs $148.3$--$152.7\,$ms).
\end{optsummary}

The full-write witness builder rewrites every byte of z/a/b each prove ---
including $b$'s all-ones prefix and reserved-slot words and all three
streams' zero tails, which are identical for every prove of a layout. A
provenance tag records which prove-layout's output a pooled buffer
currently holds. It is derived independently at the buffer's give site
(from the R1CS constants) and at its take site (from the encoder's own
constants) --- no plumbing carries it --- so a match marks the buffer as
holding exactly a previous prove's output of this layout; any other pool
custody clears the tag. On a tagged hit the builder skips the
constant-region writes. Idea from the challenge tree's scratch tokens,
re-derived; paired kill-switch A/B: $4/5$ ($147.0$--$149.0$ vs
$148.3$--$152.7$\,ms). Byte-identity is enforced by a test that compares
a tagged give/take cycle's output against a fresh, non-skipped
generation.

\paragraph{Rejected alternatives for multithreading.}
Seven mechanisms measured null or inverted on this machine, all paired:
the P-pool-only join (twice); non-temporal stores in the witness writers
(M1 ignores the hint --- third confirmation); the stripe transpose on
E-cores during the commit ($-$, bandwidth-on-bandwidth); the all-core PCS
combine; four-layer NTT fusion on aarch64 ($+19$--$26\%$, sixteen live
F128s spill); a two-block scalar witness interleave ($+40\%$ ST, GPR
blowout); a 4-wide SIMD witness builder with scalar packing (the
packing, not the hash math, is the cost); and the full lane-wise
vectorized packing network itself (\S\ref{sec:zpool} --- null once the
stripe-buffer allocation was fixed). What separates the remaining few
percent at the production shape: the challenge tree's compact round-2
anchor+delta format with symbolic round-collapsing lookaheads --- a
thousand-line-class port, priced and declined at the current bloat bar
(the campaign's implicit line-count-vs-payoff threshold) pending an
explicit call.

\section{After the main merge: the union protocol}\label{sec:union}

On 2026-08-28 the branch absorbed the repository's main line, and with
it a different statement: the \emph{union} prover (one witness for a
registry of circuit types, batch-major, count-proportional lincheck,
an assist that binds the jagged layout), the $\mathbb{F}_{2^{256}}$
two-point out-of-domain split in the Ligerito open with per-level
OOD binding and five recursive levels, and a second zerocheck --- the
AG (genus-95) univariate skip --- beside the RS one. Every entry above
was re-audited against that tree before anything new was built. Live
already: \S\ref{sec:blake3neon}, \S\ref{sec:lcfold}, the round-1 kernel
work (\S\ref{sec:pmull}, \S\ref{sec:twolane}, \S\ref{sec:zeroskip}) and
the round-1 E-core drain. Dead by protocol: \S\ref{sec:directopen} and
\S\ref{sec:directopenrefine} (the $\mathbb{F}_{2^{256}}$ split replaced
the basis open), \S\ref{sec:stripe} (the union is batch-major), and
every open entry that targeted the retired basis-combine and
sparse-prefix induce machinery. Measured against main at production
settings, the surviving wins were worth $1.027\times$ ($8/8$) --- nearly
all of it \S\ref{sec:lcfold}. What follows is the work on the new
protocol, in pipeline order as before. The benchmark shape throughout
is $m{=}32$ (a $2^{18}$-compression BLAKE3 union, $2^{25}$-word
witness), grind-free, on the same M1~Max (8P+2E); ``MT'' is the
eight-thread P-core pool.

\subsection{Round~2 on the union: the deferred-reduction port}\label{sec:r2union}

\begin{optsummary}
\item[Baseline:] Main's shared round-2 kernel accumulated the $256$-bit message partials in the four-word general-register struct that \S\ref{sec:wideneon} had replaced; the x86 arm already accumulated wide, only aarch64 had been left behind. Every register-resident primitive of \S\ref{sec:qres} was in the tree, reachable only from the dormant compact-K block.
\item[Improvement:] Route both round-2 flows (sparse and dense) through the vector-resident accumulator, then the whole fold $\to$ multiply $\to$ accumulate chain through the q-register primitives.
\item[Mechanism:] As in \S\ref{sec:wideneon} and \S\ref{sec:qres}; bit-identical (the XOR multiset is unchanged, reduction is $\Ftwo$-linear), all thirteen proof-byte pins hold.
\item[Result:] Round 2, ST: $281.7\to239.3\,$ms ($-9.2\%$, $3/3$), then $267.2\to236.1$ ($-11.6\%$, $3/3$); $-16.2\%$ cumulative, on a baseline already cheaper than the old tree's because the union's sparse path folds only the live $72\%$ of the domain.
\end{optsummary}

The compact-K rounds and the cascade tail (\S\ref{sec:compactk},
\S\ref{sec:cascade}) were wired onto the union's grinding driver in the
same pass, with the two test oracles and the transcript-identity tests
restored. They did \emph{not} keep the bench: the union's boolean region
is $72\%$ occupied, so main's support-proportional sparse route fires
first, and compact-K's producer must pay a full-domain pass to build
its tables --- $145$--$190\,$ms for rounds 2--5 plus tail against the
sparse route's $93$--$170$. The kernels are cheaper (K fold $31$--$48$
and cascade tail $11$--$15$ against the sparse tail's $44$--$103$); the
producer eats the difference. The cascade now serves dense single-run
flows, where the sparse gate fails and the 2026-08-27 certification
applies; the sparse route dispatches first everywhere else. A same-binary
route A/B settled the question in both directions at once: sparse
round 2 $43.4$ against dense $101.7$ ($8/8$), but the cascade tail
$10.6$ against the sparse tail $44.9$ ($8/8$) --- the two halves were
disjoint in opposite directions, and the win was a sparse
four-to-one lookahead kernel that neither tree had
(\S\ref{sec:foldentry}).

\subsection{The AG skip fold: a scalar twin of an existing kernel}\label{sec:agfold}

\begin{optsummary}
\item[Baseline:] The AG path's skip$\to$multilinear fold folded each $64$-bit message word through $8\times[\Fbig;256]$ tables with a scalar byte-dot: eight table loads and eight struct XORs per word. Probes put $85\%$ of the fold in these gathers and $12\%$ in the message chain, so the round-2 treatment (\S\ref{sec:r2union}) does not apply --- the accumulator is not where the time is.
\item[Improvement:] The RS row fold \code{fold\_one\_row\_neon\_q} has the identical table layout ($8\times256$ entries, contiguous, $4\,$KB stride); it is re-exported and the AG byte-dot routed through it on aarch64.
\item[Mechanism:] Same table entries, same XOR multiset --- bit-identical --- with the eight gathers issued as vector loads and XORed in the vector file.
\item[Result:] AG fold $-3.1\,$ms MT ($-7.5\%$, $3/3$), $-28\,$ms ST ($-9.0\%$, $2/2$); the first improvement to the AG path.
\end{optsummary}

Smaller than a first look at the scalar source suggested, because LLVM
already emitted reasonable code for the $16$-byte-aligned table loads;
what the kernel recovers is the remaining gap. The lesson generalizes:
when the same fold appears in two protocols, the kernel written for one
is usually a drop-in for the other, and the audit that finds it is
cheaper than a new kernel.

\subsection{The AG zerocheck as the default: dispatch, not kernels}\label{sec:agdefault}

\begin{optsummary}
\item[Baseline:] The flat prover's default zerocheck was RS; AG existed as a parallel entry with its own prove$\to$verify test, exercised only by the tower and a micro-bench. Its sparse route had been (mis)measured as its multi-thread pessimum.
\item[Improvement:] AG is the default on aarch64 (RS stays the default elsewhere, and is retained as the x86 fallback: the AG round-1 kernel is NEON-only and the CI's x86 leg proves with RS); the SHA-256 setup gets the mechanical twin of the BLAKE3 AG entry; the AG proof gets its own wire flavor so the bench reports peak memory, verify and proof size on both arms.
\item[Mechanism:] Two corrections came first. Sparse is AG's optimum on \emph{every} stage (round 1 $55.7$ vs $57.1$, fold $41.1$ vs $47.2$, tail $47.8$ vs $65.7$, zerocheck$+$lincheck $183$ vs $209\,$ms; end to end $4/4$) --- the earlier ``$3.2\times$'' reading was a timing artifact --- so the dense/sparse gate follows the flavor. Then the promotion measured in the shipped configuration, both arms on the default sparse gate.
\item[Result:] Prove median $966.1\to894.8\,$ms ($-8.2\%$, $4/4$; deltas $-112/-84/-75/-68$), best $935.0\to851.4$: worth more than every kept port of the campaign combined ($1.027\times$), and it is not new kernel code.
\end{optsummary}

The RS path cannot be deleted --- it owns every transcript-byte pin and
x86 proving --- so both entries stay, as the recursion plan prescribes.
The subsequent round-1 audit against AG found nothing above the bloat
bar: the multiplicative-weight split of \S\ref{sec:pmull} has no
reduction to defer in AG (its products are bitsliced $\Ftwo$ ANDs, and
its eq-weighting already accumulates unreduced), structured-$b$ of
\S\ref{sec:structuredb} has a $2\%$ ceiling there, and the stripe-C
derivation of \S\ref{sec:stripe} maps onto batch-major addresses
(skip $=$ in-word bits $0$--$5$, parity $=$ bit $6$, inner $=$ batch bits
$0$--$5$, outer $=$ the rest) but costs more multiplies than it removes
transposes: AG's C side is a multiply roofline and its transpose is free.

\subsection{Witness generation: attributed, and the dead regions skipped}\label{sec:witgen2}

\begin{optsummary}
\item[Baseline:] The trace printed $0.00\,$ms for witness generation because the timer wrapped a passthrough; the real generator ran outside the traced region, $64\,$ms ($7.6\%$ of the prove) that no session had attributed. Attributed: $1.5\,$GB of non-temporal stores in $25.7\,$ms (a $58\,$GB/s ``roofline''), $17.8$ of BLAKE3 compute, a $432\,$MB padding memset ($15.0$), the lincheck stripe ($12.8$) and its $153\,$MB tail clear ($7.8$), row fill ($6.2$).
\item[Improvement:] Three changes. (i) The $a/b$ padding columns and the stripe tail are never read by anything --- the run-list-gated zerocheck skips dead blocks, the lincheck kernel is handed \code{useful\_bits} --- so they are no longer cleared when the union certifies that (\code{dead\_padding\_unread}); $z$'s suffix keeps its memset. (ii) The flush stages eight consecutive groups and writes each chunk-column as one $1\,$KB burst instead of $128$-byte scatters at a $4\,$MB stride. (iii) The BLAKE3 builder's super-group loop joins the P$+$E work queue.
\item[Mechanism:] (i) is proven rather than argued: a guard test poisons both regions on one of two otherwise identical proves and pins the proofs byte-identical. The $58\,$GB/s was never the store roofline; it was the roofline for one $128$-byte store per page. (ii) keeps $96$--$384\,$KB of staging in L2 ($768\,$KB spills and gives the win back).
\item[Result:] (i) witgen $-10.9$ ($4/4$) and $-3.2\,$ms ($8/8$) on the micro-bench, $-22.4\,$ms end to end ($4/6$; more than the two skips account for --- $441\,$MB fewer dirty pages ahead of the commit). (ii) ST $351.6\to339.4$ ($-12$), MT $-3$ to $-6$. (iii) $-2.1\,$ms MT ($6/6$). In-prove witgen $64\to{\approx}35\,$ms MT.
\end{optsummary}

Two findings frame what is left. The padding memset for $z$
\emph{cannot} be skipped, and the reason is completeness, not
soundness: at full utilization the compaction is the identity, so $q$
\emph{is} the padded buffer and the committed polynomial includes the
padding while the zerocheck and lincheck claims are computed as if it
were zero --- forcing the skip makes the honest prover's own proof fail
at \code{PcsOpen(Ligerito)}, the predicted site. (Padding is written
entirely prover-side, so it can never affect what a verifier accepts.)
And \S\ref{sec:elision} was re-examined against the new generator: the
part of it that pays --- the three streams' zero tails --- is exactly
what \code{dead\_padding\_unread} now skips; the constant regions inside
the useful data are worth $\approx1\,$ms, and a provenance-tagged
witness pool was built end to end and found structurally unable to hit
(the buffers return through the pool in a different class than they are
taken from). What remains is $\approx15\,$ms of required stores at the
burst rate, $18$ of BLAKE3 compute, and $13$ for the stripe: either the
answer or the bus.

\subsection{The sparse tail's ladder and the fold-level entry}\label{sec:foldentry}

\begin{optsummary}
\item[Baseline:] The AG/sparse tail ran one round per pass over the live prefix: a fold-1 entry at $2^{26}\to2^{25}$ ($43\,$ms MT, $209$ ST) then rounds of geometrically shrinking cost. Per pass the lookahead ladder of \S\ref{sec:cascade} loses on the entry ($+16\%$) and wins on every steady pass ($-19\%$, $-21\%$).
\item[Improvement:] (i) The prefix ladder: fold-2 passes with the deferred-challenge state on the compact buffer, ``one scaling per row'' (pre-scale the four $a$ values by eq so the eight products carry the weight as single unreduced multiplies: $16\to12$ per position). (ii) The fold-level entry: the skip$\to$multilinear fold forms the round-0/1 lookahead sums itself, NEON-resident off the byte-dots of \S\ref{sec:agfold}, so the tail's first pass is already a fold-2 over the full $48$M and there is no entry pass.
\item[Mechanism:] Transcript-identical (the round-0 message falls out of the sums). The entry's cost was first priced by \emph{adding} the accumulation to the fold's existing Horner message ($+19\,$ms) and declared dead; in a fused fold the sums \emph{replace} that message, and the fair cost is the difference: $+5\,$ms with the vector-resident formation.
\item[Result:] Ladder $-4.5\,$ms MT ($2/3$); entry $-10\,$ms MT ($3/3$), $-12$ ST; tail $62\to45$--$47$.
\end{optsummary}

The mis-pricing is worth recording as a rule: when pricing a fusion,
\emph{replace, don't add}. A related null closes the subject: the
friendly-Horner tail kernel is $2$--$3\times$ slower than the ladder on
this path and is never dispatched.

\subsection{The statistics ladder in the open}\label{sec:statsladder}

\begin{optsummary}
\item[Baseline:] The open's level-0 basis is a sum of $T{=}2$ rank-one eq tensors (the merged transport's seeded point and the one level-0 OOD point), and the six initial rounds bind exactly the six lane-block bits. The prover nonetheless swept the $2^{25}$ witness for the combine's prime and lookahead ($46\,$ms MT), swept it again for the OOD ($20$), then folded it through five passes with the virtual basis evaluated per output ($63$): the same structure paid for three times.
\item[Improvement:] Write the index as $u = e\cdot d + h$; each term is $s_t\,\mathrm{blk}_t[e]\,\mathrm{within}_t[h]$. Every level-0 round message is a function of the \emph{block statistics} $G_t[e] = \sum_h f[e,h]\,\mathrm{within}_t[h]$ --- $64$ values per term --- and the six array folds compose into one $64$-term fold $f'[h] = \sum_e \eq(r,e) f[e,h]$ straight into $\mathbb{F}_{2^{256}}$. One sweep computes both terms' statistics (two multiplies per element); the OOD evaluation and $\beta$ come from the $64$-entry tables; the switch basis is $\sum_t s_t\,\mathrm{blk}_t[0]\,\mathrm{within}_t[h]$.
\item[Mechanism:] Same field elements in the same transcript slots, so proof bytes are identical; an exact oracle test at small geometry. This is the challenge tree's ``direct fold'', which the port audit had filed as dead by protocol because their \emph{code} was tied to the $100$-bit basis open. The algebra is not: it needs a rank-one basis and lane-major block rounds, which the $\mathbb{F}_{2^{256}}$ ladder has. Re-derived, not ported.
\item[Result:] Leaf open $298.9\to204.6\,$ms; in-prove open $-80$ to $-107$, prove $-115$ to $-138\,$ms ($3/3$), best prove $760.6\to630.3$. The total moves $\approx35\,$ms more than the open bucket because the ladder no longer allocates fold 1's two $256\,$MB transients and returns the witness right after its one fold: less pool churn.
\end{optsummary}

Three follow-ups each removed a pass the ladder had left. The seeded
term's statistics are the merged sumcheck's own folded witness after
$\log n - \mathrm{initial}\_k = 19$ of its $25$ rounds (the same
coordinate convention, dead blocks honest zeros), so the sweep covers
the OOD term alone: $10.4\to6.7\,$ms ($3/3$). The code switch's three
conversions (split, promote, split the basis: $80\,$MB of fresh
allocations) are gone because the fold writes its output directly in
the switch's split form and the basis is materialized already split:
switch $7.7\to0.7\,$ms ($3/3$). And the switch's message promotion had
been a \emph{serial} map over $2^{20}$ elements into a fresh allocation
beside a parallel neighbour: $9.5\to3$--$4\,$ms, with a lesson attached
--- the first version also routed the switch's buffers through the
scratch pool, which made the switch faster and the \emph{next} fold
$3\to21$--$62\,$ms: at a $24$-entry cap the most-populated-class
eviction policy threw out buffers the fold reuses. Adding pool traffic
is not a local change (\S\ref{sec:poolcap}).

\subsection{The weight build with the column factor baked into fold tables}\label{sec:bakedw}

\begin{optsummary}
\item[Baseline:] The merged transport's weight $W[d] = \sum_i \varphi_i(\eq_{\mathrm{row},i}[r]\cdot\eq_{\mathrm{col},i}[c])$ cost, per element and per RS claim, one field multiply by the hoisted column factor and then the $16$-lookup byte-table map $\varphi_i$.
\item[Improvement:] The column factor is constant across a column segment and $x\mapsto\varphi_i(\eq_{\mathrm{col}}[c]\cdot x)$ is $\Ftwo$-linear, so its byte decomposition \emph{is} a fold table: bake $\Psi_c = \varphi_i\circ(\eq_{\mathrm{col}}[c]\cdot)$ per column ($92$ live $\times\,64\,$KB per claim, $\approx1\,$ms MT, parallel over columns) and pay the sixteen lookups only.
\item[Mechanism:] The same field elements, so $W$ and everything downstream are identical (pins pass). The same principle as the virtual basis and the statistics: never materialize an eq tensor, and never multiply by a factor that a table can absorb.
\item[Result:] W build MT $32.2\to25.4\,$ms ($3/3$), ST $227\to167$.
\end{optsummary}

\subsection{Mixed-field intake for the level OODs and the induced bases}\label{sec:mixedfield}

\begin{optsummary}
\item[Baseline:] The recursive levels' OOD bases ($\eq(z,\cdot)$, $2^{20}$ at level 1) and the induced query basis ($2^{19}$, transported across the split as the pairs $(B_j,0),(0,B_j)$) are base-field valued, and so is the table they meet; the intake promoted each to $\mathbb{F}_{2^{256}}$ (a $32\,$MB allocation and write), formed the message with $\mathbb{F}_{2^{256}}\times\mathbb{F}_{2^{256}}$ products (three $\Fbig$ multiplies each), split the induced basis into a second $64\,$MB array, and glued with mixed products into both limbs. These buckets scaled only $3.5\times$ ST$\to$MT: page faults on the fresh promotions.
\item[Improvement:] A pending-basis form (\code{Base}, \code{PresplitBase}): the OOD evaluation and its $(u_0,u_2)$ are three $\Fbig$ products per pair in one sweep; the presplit message is $u_0 = \sum f_0 B_j$, $u_2 = (S,S)$ with $S = \sum (f_0{+}f_1)B_j$ (two products per pair, since $u\cdot B = (0,B)$ and $(1{+}u)s = (s,s)$ for base $s$); the glue adds $B_j\beta$ into the right limb.
\item[Mechanism:] Same field elements; no promotion, no split array; the promoted intake deleted.
\item[Result:] Level OODs $10.4\to4.6$--$5.6\,$ms MT ($35.9\to11.7$ ST), induce $10$--$11\to7$ MT ($36.2\to23$ ST); open $-16/-17$; best prove $601.4\,$ms.
\end{optsummary}

\subsection{The level-0 Merkle was hashing every leaf on the generic path}\label{sec:leafbatch}

\begin{optsummary}
\item[Baseline:] \S\ref{sec:blake3neon}'s eight-message kernel was dispatched only for leaf sizes $64/128/256/512/1024$ (``leaf sizes are $16\ll\log$ batch, so only powers of two arise'' --- true before the integer-lane commit). The union commits $46$ lanes at $m{=}32$: a $736$-byte leaf, so every one of the $2^{20}$ leaves went to the generic single-stream hasher. The commit attribution had called this ``at the hash roofline for rate $1/4$'' on a $2\,$GB codeword; the level-0 rate is $1/2$ and the stored codeword is $736\,$MiB, so the Merkle was running at $6.3\,$GB/s --- a third of the kernel's rate.
\item[Improvement:] The kernel takes the last block's true length (a short final block is zero-padded, as BLAKE3 specifies) and the batcher accepts any single-chunk leaf ($1$--$1024$ bytes), tail leaves on the generic path.
\item[Mechanism:] Same chunk chaining values, same roots (pins pass; $24$ unit tests over the new sizes, $736$ among them). The misattribution is the lesson: a ``roofline'' claim has to be checked against the actual byte count.
\item[Result:] Merkle MT $116.7\to39.8\,$ms ($3/3$), $19.8\,$GB/s on the all-core pool (the challenge tree's rate is $21.6$); commit $232.9\to154.4$; prove $-67$ to $-90$; ST Merkle $941\to314$, commit $1640\to1008$.
\end{optsummary}

This fix answered a question the numbers had been hiding: with the AG
skip so much cheaper than RS round 1, why had commit $+$ zerocheck $+$
lincheck not fallen? Before the fix the sum was $405$--$415\,$ms, the same
as the pre-merge $401$ --- AG's saving (RS round 1 $165\,$ms all-in
$\to$ AG $48$) had been cancelled by round 1 no longer having an idle
commit window to hide under, and by the Merkle at a third of its rate.
After it the sum is $\approx340$.

\subsection{Block-first merged sumcheck from column statistics}\label{sec:blockfirst}

\begin{optsummary}
\item[Baseline:] The merged transport (the ring-switching note, \emph{docs/capacity-free-ring-switching.tex}) materialized its weight ($25\,$ms) and ran a $25$-round product sumcheck over the $2^{25}$ domain ($46\,$ms) before the open.
\item[Improvement:] Under a full-height column prefix the weight is a sum of rank-one terms ($256$ per RS claim: a column bit-bank $T[c][b] = \sum_r \eq(z_{\mathrm{row}},r)\,\mathrm{bit}_b(q[c,r])$ against a column factor), so the block coordinates are bound first from block statistics obtained by a Frobenius transform of the bit banks, and only the last rounds touch the rows.
\item[Mechanism:] Every weight term masked to the live prefix (the verifier's $\widehat W$ is the zero-tail extension); $\rho$ rotated into coordinate order in prover, verifier and the recursive walker. The C claim's bit bank is one extra fold of the lincheck stripe at the C row point ($13.6\,$ms, outside the open).
\item[Result:] Merged sumcheck $46\to14\,$ms, weight build gone; open $137\to109\,$ms MT, prove $\approx-15\,$ms net of the C fold ($3/3$).
\end{optsummary}

The full fusion (the weight as the open's level-0 basis, no merged
sumcheck at all) was built to the native verifier and parked: the
residual side of the recursive verifier does not twist for free, and
its $+8.5$k rows per chain child buy a one-percent leaf gain. Fusing the
C fold into the lincheck's own sweep (a two-point stripe kernel with
$32$-byte sum-table entries and sixteen accumulators) was built,
bit-identical, and measured a wash ($-0.3\,$ms): the stripe kernel is
core-throughput-bound at $\approx1.4$ lookups per cycle per core
($85.9\,$ms on one thread, $14.4$ on eight, $37\,$GB/s), not
bandwidth-bound, so a second point costs its own load and XOR per
stripe byte and shares only the addressing. With the $q'$ fold and the
ladder's seed sweep at the multiply rate, stage 1's remaining
$\approx13\,$ms of transport tax is protocol-shaped.

\subsection{Lane-fill tiling and first-prove latency}\label{sec:lanefill}

\begin{optsummary}
\item[Baseline:] The commit's lane-major fill tiled at $64$ positions, and the tile length is the per-lane read burst: $1\,$KB, with $t$ lanes whose sources sit $8\,$MiB apart. The prefetcher never establishes a stream across $61$ far-separated $1\,$KB bursts; a pure copy ran at $9\,$GB/s.
\item[Improvement:] A size-aware tile, $\mathrm{tile} = (d/(4\,\mathrm{threads}))$ clamped to $[64, 1024]$ positions, keeping at least four tiles per worker.
\item[Mechanism:] Longer bursts per lane; the same words to the same addresses.
\item[Result:] ST fill $136\to94\,$ms ($-31\%$; $1024$ positions in the sweep); MT a wash --- at eight workers the pass is bandwidth-saturated at any burst length.
\end{optsummary}

Two latency items, invisible to a warm minimum and important to a
one-shot prover: the statement-binding digest (a BLAKE3 hash of every
type's sparse matrices, $\approx21$M nonzeros) was materialized inside
the \emph{first} prove, $1.3\,$s; it now joins the setup's warm-up
(\code{bind statement} $1297\to0.01\,$ms on the first prove). And the
first open in a process is $\approx100\,$ms colder than the warm one
($279$ against $181$ at the time): a pool eviction that unmaps a
multi-hundred-megabyte buffer, plus cold recursive folds.

\subsection{Scratch-pool retention: the count was the constraint}\label{sec:poolcap}

\begin{optsummary}
\item[Baseline:] The scratch pool retained at most $24$ buffers, evicting from the most-populated size class. At steady state the $m{=}32$ prove cycles $\approx25$ (a dozen small ladder and query buffers plus the witness $\times3$, the codeword, the zerocheck's four $2^{25.5}$-word giants and the fold outputs), so the pool overflowed on every prove and evicted the giants during the open's gives; the next zerocheck missed four times and re-mapped $\approx3\,$GB of fresh pages --- $\approx197$k page reclaims per prove.
\item[Improvement:] Retention $24\to48$.
\item[Mechanism:] Peak RSS is unchanged ($7.2$--$7.9\,$GB in every configuration --- the buffers exist during every prove anyway; retention only adds idle footprint). A warm-first (LIFO) tie-break was tried in the same experiment and is harmful in both settings; the 2026-07-28 byte-budget experiment recorded in the source had moved nothing because it ranked eviction by bytes at the same \emph{count} cap.
\item[Result:] Same binary, eight-prove steady state, four alternating pairs: prove $-55$ to $-96\,$ms (mean $-68$, $4/4$), zerocheck$+$lincheck $-39$, misses $45\to12$ per nine proves, page reclaims $2.0$M$\to1.0$M; best prove $444.5$, then $438.3\,$ms on a quiet box.
\end{optsummary}

Two follow-ups closed the subject. There is no other leak: steady-state
faults are $660$ per prove and RSS is flat over eight proves; the
fault storms this desktop shows under load are the memory compressor
reclaiming the retired witness's pages (the ``witgen cold-page
lottery'': the same phase reads $3$--$4\,$ms with its pages resident and
$44$--$100$ when they were compressed, which is why an A/B compares only
the phase it changed). And wiring the pool's large buffers with
\code{mlock} measured nothing on a quiet box ($4$ pairs, mixed signs
within $\pm10\,$ms on min, median and max) while adding $\approx63$k
faults per prove for the wiring walks.

\section*{Summary}

\paragraph{Measurement protocol.} Results before 2026-08-27 include
proof-of-work grinding; later entries marked \emph{grind-free} zero the
grinding bits on both trees (\code{FLOCK\_NO\_GRIND=1}; the library
verifiers stay honest, so such proofs fail verification by design) —
nonce-search luck otherwise adds variance paired runs cannot cancel.
Grind-free figures form their own baseline family. Sub-3-ms phase-local
effects are certified by pairing the phase line itself rather than
bucket or total times. From Section~\ref{sec:union} on, the shape is
the $m{=}32$ union at production settings, measured same-binary where a
knob allows it, otherwise alternating saved binaries; steady-state
figures take the minimum over eight proves in one process, and a phase
whose pages the memory compressor may have reclaimed (witness
generation, on this desktop) is compared only when it is the phase
under test.


\begin{table}[h]
\centering
\small
\begin{tabular}{@{}llrr@{}}
\toprule
Section & Kind & Phase effect (ST) & Sign test\\
\midrule
\S\ref{sec:wideneon} & representation & round 2: $-13.4\%$ & decisive\\
\S\ref{sec:zeroskip} & work removal & round 1: $-2.9\%$ & $8/8$\\
\S\ref{sec:twolane} & ILP & round 1: $-4.3\%$ & $8/8$\\
\S\ref{sec:pmull} & arithmetic & round 1: $-9.1\%$ & $8/8$\\
\S\ref{sec:stripe} & structural & round 1: $-5.7\%$ & $8/8$\\
\S\ref{sec:qres} & representation & round 2: $-13.0\%$ & $8/8$\\
\S\ref{sec:tailfuse} & pass fusion & rounds 3+: $-31.8\%$ & $8/8$\\
\S\ref{sec:composedfold} & work removal & open combine: $-17\%$ & decisive\\
\S\ref{sec:directopen} & work removal & open ($m{=}32$): $-34\%$ 8T, $-66\%$ ST & $3/3$ disjoint\\
\S\ref{sec:directopenrefine} & pass fusion & open ($m{=}32$): $-20\%$ 8T & $3/3$ disjoint\\
\S\ref{sec:cascade} & pass fusion & zc tail ($m{=}32$): $-38\%$ 8T & $4/4$ disjoint\\
\S\ref{sec:compactk} & representation, pass fusion & zc r2..5 ($m{=}32$): $-4.9\,$ms 8T & $3/3$ disjoint\\
\S\ref{sec:structuredb} & work removal, scheduling & commit join / zc r1 ($m{=}32$): $-6.2\,$ms 8T & $5/6$, $3/3$\\
\S\ref{sec:inducetrunc}+\S\ref{sec:lazyood} & work removal & open ($m{=}32$): $-9\%$ 8T & $3/3$ disjoint\\
\S\ref{sec:recursivefill} & pass fusion & open ($m{=}32$): $-1.5\,$ms 8T & sign $3/3$\\
\S\ref{sec:boundaryjoin}+\S\ref{sec:blockedtranspose} & scheduling, pass fusion & open ($m{=}32$): $-12\%$ 8T & sign $3/3$\\
\S\ref{sec:streamcommit} & architecture & commit+prep ($m{=}32$): $-3\%$ wall, $-5\%$ total & $3/3$ disjoint\\
\S\ref{sec:rankedtop} & kernel, scheduling & commit+prep ($m{=}32$): $-5\%$ wall & $3/3$ disjoint\\
\S\ref{sec:streamwit} & work removal & witness gen: $-24\%$ & $3/3$\\
\S\ref{sec:lcfold} & work removal & lincheck fold: $-32\%$ & $3/3$\\
\S\ref{sec:blake3neon} & ILP & BLAKE3 merkle: $+41\%$ GB/s & decisive\\
\midrule
\multicolumn{4}{@{}l}{\emph{on the union protocol, $m{=}32$ ($2^{18}$ compressions), grind-free}}\\
\S\ref{sec:r2union} & representation & round 2: $-16.2\%$ ST & $3/3$, $3/3$\\
\S\ref{sec:agfold} & kernel reuse & AG fold: $-7.5\%$ 8T, $-9.0\%$ ST & $3/3$, $2/2$\\
\S\ref{sec:agdefault} & dispatch & prove: $-8.2\%$ 8T & $4/4$\\
\S\ref{sec:witgen2} & work removal, store pattern & witness gen: $64\to{\approx}35\,$ms 8T; prove $-22\,$ms & $4/6$, $8/8$, $6/6$\\
\S\ref{sec:foldentry} & pass fusion & zc tail: $-14.5\,$ms 8T, $-12$ ST & $2/3$, $3/3$\\
\S\ref{sec:statsladder} & work removal (algebra) & open: $-80$ to $-107\,$ms 8T; prove $-115$ to $-138$ & $3/3$\\
\S\ref{sec:bakedw} & representation & W build: $-7\,$ms 8T, $-60$ ST & $3/3$\\
\S\ref{sec:mixedfield} & representation & level OODs $+$ induce: $-9\,$ms 8T, $-37$ ST & $3/3$\\
\S\ref{sec:leafbatch} & dispatch & merkle: $116\to39\,$ms 8T; commit $-78$ & $3/3$\\
\S\ref{sec:blockfirst} & protocol transport & open: $137\to109\,$ms 8T & $3/3$\\
\S\ref{sec:lanefill} & store pattern & lane fill: $-31\%$ ST & sweep\\
\S\ref{sec:poolcap} & allocation & prove: $-68\,$ms 8T (mean) & $4/4$\\
\midrule
\multicolumn{2}{@{}l}{round 1 cumulative} &
  \multicolumn{2}{r@{}}{$1477 \to 1205\,$ms\quad($-18.4\%$)}\\
\multicolumn{2}{@{}l}{round 2 cumulative
  (\S\ref{sec:wideneon}+\S\ref{sec:qres})} &
  \multicolumn{2}{r@{}}{$433 \to 312\,$ms\quad($-28\%$)}\\
\multicolumn{2}{@{}l}{rounds 3+ cumulative} &
  \multicolumn{2}{r@{}}{$455 \to 307\,$ms\quad($-33\%$)}\\
\multicolumn{2}{@{}l}{PCS open cumulative (\S\ref{sec:composedfold}+\S\ref{sec:directopen}, $m{=}30$ ST / $m{=}32$ ST)} &
  \multicolumn{2}{r@{}}{$\approx160 \to \approx145\,$ms\;/\;$614\to210\,$ms}\\
\multicolumn{2}{@{}l}{end-to-end, all seven zerocheck changes} &
  \multicolumn{2}{r@{}}{$\approx52{,}400 \to \approx58{,}100$ comp/s\quad($\approx+11\%$, $8/8$)}\\
\multicolumn{2}{@{}l}{whole campaign on the old protocol, $m{=}32$ BLAKE3 (grind-free, certified final table)} &
  \multicolumn{2}{r@{}}{8T $723.9 \to 454.7\,$ms ($1.59\times$); ST $4781 \to 3230$ ($1.48\times$); gap $1.36\times\to1.217\times$ MT / $1.145\times$ ST}\\
\multicolumn{2}{@{}l}{union protocol, $m{=}32$ BLAKE3 (grind-free), from the merge to now} &
  \multicolumn{2}{r@{}}{8T $853.8 \to 438.0\,$ms ($1.95\times$); ST $5.25 \to 2.91\,$s ($1.80\times$); the retired direct pipeline's $458.9$ / $3.28$ beaten on both}\\
\bottomrule
\end{tabular}
\end{table}

\paragraph{Rejected alternatives.}
The kept changes are one side of a 17-experiment record; the failures shaped
them. On this core, re-encoding fixed work never paid: pairing multiplies
through a memory-interfaced kernel ($+10\%$), a four-register structured load
($+12\%$, microcoded), three-way XOR fusion ($+1.8\%$), an $8\times$ smaller
gather table ($-3.6\%$ regression), and deferring an already-cheap reduction
($0\%$) all failed, while every win either removed work, added independent
work in flight, or removed register-file boundary crossings. In the
commit-phase NTT, rewriting the butterfly multiply as a $3$-PMULL Karatsuba
with a q-resident interface regressed $58\%$: the generic butterfly already
inlines the latency-optimal binius kernel, and six PMULLs with no repack beat
three with one, in situ and not only in the microbenchmark. Two further
structural transplants measured neutral end-to-end and were reverted despite
making the zerocheck \emph{bucket} look better (hoisting the challenge-independent
prep under the commit; the lookahead double-round --- which first lost $7$--$15\%$ in the
zerocheck tail and was rejected here, until the cross-tree kernel anatomy
showed the loss belonged to the scalar accumulator, not the technique:
re-kerneled it is now \S\ref{sec:cascade}). The commit-phase campaign
added measured nulls of its own: pairing the AB prep with the Merkle
stage on the theory that BLAKE3 (integer SIMD) and PMULL use disjoint
ports (both issue on the NEON pipes; the Merkle arm went $55\to150\,$ms
beside the prep); hashing leaves \emph{inside} the NTT's deep tasks
(serializes BLAKE3 behind PMULL on the same core — the win required the
E-core \emph{pipeline} of \S\ref{sec:streamcommit}); a 16-stream NEON
fused-4 top ($2.5\times$ worse: prefetch breakdown at huge strides) and
a first fused-3 whose per-lane scatter stores paid read-for-ownership on
the fresh destination — radix-8 only won as \S\ref{sec:rankedtop}'s
staged-store kernel, a reminder that a null adjudicates one
implementation in one architecture, not the idea. A pooled induce
densify buffer also measured null (the right-sizing copy-out ate the
fault savings at 32\,MB scale), and the whole-prove E-core pass kept
only one of three sites: hetero tiles on the open's boundary sweep and
recursive-commit leaves measured negative --- on this host the pattern
needs $\gtrsim$30\,ms compute-bound windows to amortize its helper
threads.

\paragraph{Rejected alternatives, after the merge.}
Section~\ref{sec:union}'s twelve kept changes sit beside about twenty
measured refusals, and several of them are the more instructive
record. \emph{Scheduling nulls}: hoisting round-1 AB prep under the
union's commit moved its buckets exactly as designed (round 1
$120\to55$) and the commit paid it all back ($+52$, $0/6$) --- the NTT
is at $50\,$GB/s with its top tiles already on P$+$E and the Merkle at
the hash roofline on the all-core pool, so the window has no idle
capacity to hoist into; the same holds for AG round 1, which is more
hoistable in principle. The lane-major streaming commit measured
$6/12$ at $+2\,$ms, i.e.\ zero. The E-core queue that pays on round 1
and the witness builder loses $+24\,$ms ($3/3$) on the skip$\to$mlv
fold and $+3$ on the statistics sweep: E-cores help only flat,
compute-bound loops, never a phase at the bus. \emph{Kernel re-encodings}:
$16$-bit gather tables ($4$ gathers from $4\,$MiB instead of $8$ from
$32\,$KiB) regressed $21$--$25\%$; a Karatsuba accumulate for AG round 1
has no reduction to defer; structured-$b$ for AG has a $2\%$ ceiling;
stripe-C for AG derives cleanly onto batch-major addresses and costs
more multiplies than it removes transposes; two transposed-NTT layers
per pass in the induced-basis builder were worth $0.3\,$ms
($16\,$MB, SLC-resident); a tiled NEON radix-8 in the recursive
commits' NTT gained $0.5$ on the top pass and lost $1.0$ on the deep
one. \emph{Fusions priced honestly}: an alternating skip/double-fold
schedule on the merged product sumcheck (two variants, both
byte-identical) saved traffic on paper and lost on the wall
($\mathrm{W}$ build $+5$, or ST $+26$); having witness generation
compute round 1's challenge-free half is negative because round 1's
reads are not exposed ($38.5$ of its $46\,$ms is compute) and the
intermediate would cost more to re-read than the $8\,$ms it saves;
interleaving the witness builder's stores with its compute was built,
bit-identical, and measured nothing; the two-point stripe fold
(\S\ref{sec:blockfirst}) shares only the addressing. \emph{Structural}: a
provenance-tagged witness pool cannot hit; a four-level Fast/Slim
ladder saves $12.9\,$KB of proof natively but the recursive verifier's
row envelope cannot hold the $16\times$ residual; the full fusion of the
transport into the open is parked for the same verifier. \emph{Memory}:
\code{mlock} on the pool, and a LIFO pool tie-break, both refuted. The
one recurring mistake was pricing a fusion by adding its new work to
the old (\S\ref{sec:foldentry}); the one recurring success was reading
the cost accounting before writing a kernel (\S\ref{sec:leafbatch},
\S\ref{sec:witgen2}, \S\ref{sec:poolcap}).

\paragraph{Where the pipeline stands.} Three reference points, all on
this machine, $m{=}32$, grind-free (multi-thread / single-thread): the
retired direct pipeline at the branch's pre-merge tip, $458.9\,$ms /
$3.28\,$s with a $437\,$KB proof; the union pipeline as merged,
$853.8$ / $5.25$; the union pipeline now, $438.0$ / $2.91$ with a
$462\,$KB proof. Against the direct pipeline the commit is $-42\%$
(\S\ref{sec:leafbatch}, \S\ref{sec:blake3neon}), zerocheck $+$ lincheck
$+16\%$ (a faster zerocheck, but the union lincheck and the C-claim
bank are new work), and the open $3.6\times$ ($105$ against $30\,$ms:
two ring-switched claims and the merged transport where the old path
had one direct open --- protocol, not implementation). The independent
challenge tree (the Yukon \emph{flock-challenge} frontier: $\Fbig$
statement, SHA-256 Merkle, $19$-bit fold grinding, no union, plus a
Metal GPU commit and zerocheck fold) proves the same $2^{18}$
compressions in $266\,$ms here as the benchmark runs it, $384$--$407$
with its GPU switched off, and $2.86\,$s single-threaded: the CPU work
is equal, its CPU-only edge is parallel efficiency (an E-core helper
pool and keep-alive spinners), and its real edge is the GPU.

\end{document}
